\documentclass[11pt]{article}

% ---- theme variables (passed via pandoc -V) ----
\newcommand{\ThemeName}{Algorithms \& Complexity}
\newcommand{\ThemeKicker}{DSA Track}
\newcommand{\ThemeNumber}{03}

\usepackage{interviewstyle}
\setthemecolor{be123c}{fb7185}

% pandoc-required plumbing
\providecommand{\tightlist}{\setlength{\itemsep}{1pt}\setlength{\parskip}{0pt}}
\setlength{\emergencystretch}{3em}
\providecommand{\pandocbounded}[1]{#1}

% syntax-highlighting macros emitted by pandoc (skylighting)
\usepackage{color}
\usepackage{fancyvrb}
\newcommand{\VerbBar}{|}
\newcommand{\VERB}{\Verb[commandchars=\\\{\}]}
\DefineVerbatimEnvironment{Highlighting}{Verbatim}{commandchars=\\\{\}}
% Add ',fontsize=\small' for more characters per line
\usepackage{framed}
\definecolor{shadecolor}{RGB}{35,38,41}
\newenvironment{Shaded}{\begin{snugshade}}{\end{snugshade}}
\newcommand{\AlertTok}[1]{\textcolor[rgb]{0.58,0.85,0.30}{\textbf{\colorbox[rgb]{0.30,0.12,0.14}{#1}}}}
\newcommand{\AnnotationTok}[1]{\textcolor[rgb]{0.25,0.50,0.35}{#1}}
\newcommand{\AttributeTok}[1]{\textcolor[rgb]{0.16,0.50,0.73}{#1}}
\newcommand{\BaseNTok}[1]{\textcolor[rgb]{0.96,0.45,0.00}{#1}}
\newcommand{\BuiltInTok}[1]{\textcolor[rgb]{0.50,0.55,0.55}{#1}}
\newcommand{\CharTok}[1]{\textcolor[rgb]{0.24,0.68,0.91}{#1}}
\newcommand{\CommentTok}[1]{\textcolor[rgb]{0.48,0.49,0.49}{#1}}
\newcommand{\CommentVarTok}[1]{\textcolor[rgb]{0.50,0.55,0.55}{#1}}
\newcommand{\ConstantTok}[1]{\textcolor[rgb]{0.15,0.68,0.68}{\textbf{#1}}}
\newcommand{\ControlFlowTok}[1]{\textcolor[rgb]{0.99,0.74,0.29}{\textbf{#1}}}
\newcommand{\DataTypeTok}[1]{\textcolor[rgb]{0.16,0.50,0.73}{#1}}
\newcommand{\DecValTok}[1]{\textcolor[rgb]{0.96,0.45,0.00}{#1}}
\newcommand{\DocumentationTok}[1]{\textcolor[rgb]{0.64,0.20,0.25}{#1}}
\newcommand{\ErrorTok}[1]{\textcolor[rgb]{0.85,0.27,0.33}{\underline{#1}}}
\newcommand{\ExtensionTok}[1]{\textcolor[rgb]{0.00,0.60,1.00}{\textbf{#1}}}
\newcommand{\FloatTok}[1]{\textcolor[rgb]{0.96,0.45,0.00}{#1}}
\newcommand{\FunctionTok}[1]{\textcolor[rgb]{0.56,0.27,0.68}{#1}}
\newcommand{\ImportTok}[1]{\textcolor[rgb]{0.15,0.68,0.38}{#1}}
\newcommand{\InformationTok}[1]{\textcolor[rgb]{0.77,0.36,0.00}{#1}}
\newcommand{\KeywordTok}[1]{\textcolor[rgb]{0.81,0.81,0.76}{\textbf{#1}}}
\newcommand{\NormalTok}[1]{\textcolor[rgb]{0.81,0.81,0.76}{#1}}
\newcommand{\OperatorTok}[1]{\textcolor[rgb]{0.81,0.81,0.76}{#1}}
\newcommand{\OtherTok}[1]{\textcolor[rgb]{0.15,0.68,0.38}{#1}}
\newcommand{\PreprocessorTok}[1]{\textcolor[rgb]{0.15,0.68,0.38}{#1}}
\newcommand{\RegionMarkerTok}[1]{\textcolor[rgb]{0.16,0.50,0.73}{\colorbox[rgb]{0.08,0.19,0.26}{#1}}}
\newcommand{\SpecialCharTok}[1]{\textcolor[rgb]{0.24,0.68,0.91}{#1}}
\newcommand{\SpecialStringTok}[1]{\textcolor[rgb]{0.85,0.27,0.33}{#1}}
\newcommand{\StringTok}[1]{\textcolor[rgb]{0.96,0.31,0.31}{#1}}
\newcommand{\VariableTok}[1]{\textcolor[rgb]{0.15,0.68,0.68}{#1}}
\newcommand{\VerbatimStringTok}[1]{\textcolor[rgb]{0.85,0.27,0.33}{#1}}
\newcommand{\WarningTok}[1]{\textcolor[rgb]{0.85,0.27,0.33}{#1}}

\begin{document}

\makecoverpage

\thispagestyle{empty}
{\hypersetup{hidelinks}\tableofcontents}
\clearpage

\begin{quote}
Interview-ready reference: first principles, Python snippets, comparison tables, and cheat sheets.
\end{quote}

\section{Overview}\label{overview}

Algorithmic thinking is the foundation of every coding interview. You are not just expected to solve problems --- you are expected to \textbf{reason about efficiency}, identify optimal approaches, and communicate trade-offs clearly. This file covers the complete algorithmic toolkit needed for top-tier software engineering interviews.

\textbf{What interviewers evaluate:}

\begin{enumerate}
\def\labelenumi{\arabic{enumi}.}
\tightlist
\item
  Can you identify what kind of problem this is? (sorting, graph, DP, greedy\ldots)
\item
  Do you know the right algorithm and its complexity?
\item
  Can you implement it correctly under pressure?
\item
  Can you articulate trade-offs and edge cases?
\end{enumerate}

\textbf{Core skill loop:} Recognize \(\rightarrow\) Design \(\rightarrow\) Analyze \(\rightarrow\) Implement \(\rightarrow\) Verify \(\rightarrow\) Optimize

\section{Complexity Analysis}\label{complexity-analysis}

\subsection{Big-O / Big-Theta / Big-Omega}\label{big-o--big-theta--big-omega}

Complexity notation describes \textbf{how resource usage scales with input size}, not actual runtime.

\begin{longtable}[]{@{}
  >{\raggedright\arraybackslash}p{(\columnwidth - 4\tabcolsep) * \real{0.1429}}
  >{\raggedright\arraybackslash}p{(\columnwidth - 4\tabcolsep) * \real{0.2738}}
  >{\raggedright\arraybackslash}p{(\columnwidth - 4\tabcolsep) * \real{0.5833}}@{}}
\toprule\noalign{}
\rowcolor{acc!16}
\begin{minipage}[b]{\linewidth}\raggedright
Notation
\end{minipage} & \begin{minipage}[b]{\linewidth}\raggedright
Name
\end{minipage} & \begin{minipage}[b]{\linewidth}\raggedright
Meaning
\end{minipage} \\
\midrule\noalign{}
\endhead
\bottomrule\noalign{}
\endlastfoot
O(f(n)) & Big-O (upper bound) & Algorithm runs in \textbf{at most} f(n) steps (worst case) \\
\rowcolor{zebra}
\(\Omega\)(f(n)) & Big-Omega (lower bound) & Algorithm runs in \textbf{at least} f(n) steps (best case) \\
\(\Theta\)(f(n)) & Big-Theta (tight bound) & Algorithm runs in \textbf{exactly} f(n) steps (avg case matches both) \\
\end{longtable}

\textbf{Formal definition --- Big-O:} T(n) = O(f(n)) if \(\exists\) constants c \textgreater{} 0 and n\(_0\) such that T(n) \(\leq\) c\(\cdot\)f(n) for all n \(\geq\) n\(_0\).

\textbf{Complexity hierarchy (slowest to fastest growth):}

\setcodelang{}

\begin{verbatim}
O(1) < O(log n) < O(√n) < O(n) < O(n log n) < O(n²) < O(n³) < O(2ⁿ) < O(n!)
\end{verbatim}

\textbf{Dropping constants and lower-order terms:}

\begin{itemize}
\tightlist
\item
  O(3n\(^2\) + 5n + 12) \(\rightarrow\) O(n\(^2\))
\item
  O(100) \(\rightarrow\) O(1)
\item
  O(n + n log n) \(\rightarrow\) O(n log n)
\end{itemize}

\textbf{Multi-variable complexity:} When input has multiple dimensions, keep all:

\begin{itemize}
\tightlist
\item
  O(V + E) for graph algorithms (V = vertices, E = edges)
\item
  O(m\(\cdot\)n) for 2D DP problems
\end{itemize}

\subsection{Amortized Analysis}\label{amortized-analysis}

Amortized analysis gives the \textbf{average cost per operation over a sequence}, even if individual operations are expensive.

\textbf{Example --- Dynamic Array (Python list append):}

\begin{itemize}
\tightlist
\item
  Most appends: O(1)
\item
  Occasional doubling: O(n)
\item
  Amortized cost per append: O(1) --- because doubling happens exponentially rarely
\end{itemize}

\textbf{Three methods:}
\textbar{} Method \textbar{} Idea \textbar{} When to use \textbar{}
\textbar-\/-\/-\/-\/-\/-\/-\/-\textbar-\/-\/-\/-\/-\/-\textbar-\/-\/-\/-\/-\/-\/-\/-\/-\/-\/-\/-\/-\textbar{}
\textbar{} Aggregate \textbar{} Total cost / n operations \textbar{} Simplest; works when all operations have same amortized cost \textbar{}
\textbar{} Accounting \textbar{} Prepay credits for future expensive ops \textbar{} When operations have different costs \textbar{}
\textbar{} Potential \textbar{} \(\Phi\)(state) function tracks "stored energy" \textbar{} Most general; used for complex data structures \textbar{}

\textbf{Key amortized results:}

\begin{itemize}
\tightlist
\item
  Dynamic array append: O(1) amortized
\item
  Union-Find with union-by-rank + path compression: O(\(\alpha\)(n)) amortized \(\approx\) O(1)
\item
  Splay tree operations: O(log n) amortized
\item
  Fibonacci heap decrease-key: O(1) amortized
\end{itemize}

\subsection{Computing Time \& Space Complexity}\label{computing-time--space-complexity}

\textbf{Time complexity rules:}

\begin{enumerate}
\def\labelenumi{\arabic{enumi}.}
\tightlist
\item
  \textbf{Sequential statements:} Add costs --- O(A) + O(B) = O(max(A, B))
\item
  \textbf{Nested loops:} Multiply --- outer \(\times\) inner
\item
  \textbf{Function calls:} Substitute the called function\textquotesingle s complexity
\item
  \textbf{Recursion:} Use recurrence relations or Master Theorem
\end{enumerate}

\setcodelang{Python}

\begin{Shaded}
\begin{Highlighting}[]
\CommentTok{\# O(n) — single loop}
\ControlFlowTok{for}\NormalTok{ i }\KeywordTok{in} \BuiltInTok{range}\NormalTok{(n):}
\NormalTok{    do\_constant\_work()}

\CommentTok{\# O(n²) — nested loops}
\ControlFlowTok{for}\NormalTok{ i }\KeywordTok{in} \BuiltInTok{range}\NormalTok{(n):}
    \ControlFlowTok{for}\NormalTok{ j }\KeywordTok{in} \BuiltInTok{range}\NormalTok{(n):}
\NormalTok{        do\_constant\_work()}

\CommentTok{\# O(n log n) — loop + halving inner}
\ControlFlowTok{for}\NormalTok{ i }\KeywordTok{in} \BuiltInTok{range}\NormalTok{(n):}
\NormalTok{    j }\OperatorTok{=}\NormalTok{ n}
    \ControlFlowTok{while}\NormalTok{ j }\OperatorTok{\textgreater{}} \DecValTok{1}\NormalTok{:}
\NormalTok{        j }\OperatorTok{//=} \DecValTok{2}   \CommentTok{\# log n iterations}

\CommentTok{\# O(log n) — binary search}
\NormalTok{lo, hi }\OperatorTok{=} \DecValTok{0}\NormalTok{, n}
\ControlFlowTok{while}\NormalTok{ lo }\OperatorTok{\textless{}}\NormalTok{ hi:}
\NormalTok{    mid }\OperatorTok{=}\NormalTok{ (lo }\OperatorTok{+}\NormalTok{ hi) }\OperatorTok{//} \DecValTok{2}
    \CommentTok{\# halves search space each time}
\end{Highlighting}
\end{Shaded}

\textbf{Space complexity --- what to count:}

\begin{itemize}
\tightlist
\item
  Input space (usually excluded unless asked for auxiliary space)
\item
  Call stack depth for recursion
\item
  Extra data structures allocated
\item
  \textbf{Auxiliary space} = space excluding input
\end{itemize}

\setcodelang{Python}

\begin{Shaded}
\begin{Highlighting}[]
\CommentTok{\# O(1) space — no extra allocation}
\KeywordTok{def}\NormalTok{ find\_max(arr):}
\NormalTok{    m }\OperatorTok{=}\NormalTok{ arr[}\DecValTok{0}\NormalTok{]}
    \ControlFlowTok{for}\NormalTok{ x }\KeywordTok{in}\NormalTok{ arr:}
\NormalTok{        m }\OperatorTok{=} \BuiltInTok{max}\NormalTok{(m, x)}
    \ControlFlowTok{return}\NormalTok{ m}

\CommentTok{\# O(n) space — stores extra array}
\KeywordTok{def}\NormalTok{ prefix\_sum(arr):}
\NormalTok{    ps }\OperatorTok{=}\NormalTok{ [}\DecValTok{0}\NormalTok{] }\OperatorTok{*}\NormalTok{ (}\BuiltInTok{len}\NormalTok{(arr) }\OperatorTok{+} \DecValTok{1}\NormalTok{)  }\CommentTok{\# O(n)}
    \ControlFlowTok{for}\NormalTok{ i, x }\KeywordTok{in} \BuiltInTok{enumerate}\NormalTok{(arr):}
\NormalTok{        ps[i}\OperatorTok{+}\DecValTok{1}\NormalTok{] }\OperatorTok{=}\NormalTok{ ps[i] }\OperatorTok{+}\NormalTok{ x}
    \ControlFlowTok{return}\NormalTok{ ps}

\CommentTok{\# O(n) space — call stack for recursion}
\KeywordTok{def}\NormalTok{ factorial(n):}
    \ControlFlowTok{if}\NormalTok{ n }\OperatorTok{\textless{}=} \DecValTok{1}\NormalTok{: }\ControlFlowTok{return} \DecValTok{1}
    \ControlFlowTok{return}\NormalTok{ n }\OperatorTok{*}\NormalTok{ factorial(n }\OperatorTok{{-}} \DecValTok{1}\NormalTok{)  }\CommentTok{\# n frames on stack}
\end{Highlighting}
\end{Shaded}

\subsection{Recursion Tree \& Master Theorem}\label{recursion-tree--master-theorem}

\textbf{Recursion tree method:} Draw the tree of recursive calls, compute cost at each level, sum all levels.

\setcodelang{}

\begin{verbatim}
T(n) = 2T(n/2) + n    (merge sort)

Level 0:  n                    → cost n
Level 1:  n/2 + n/2            → cost n
Level 2:  n/4 + n/4 + n/4 + n/4 → cost n
...
Level log n: n leaves of size 1 → cost n

Total: n × (log n + 1) = O(n log n)
\end{verbatim}

\textbf{Master Theorem:} For recurrences of the form T(n) = a\(\cdot\)T(n/b) + f(n) where a \(\geq\) 1, b \textgreater{} 1:

Let critical exponent: \emph{\emph{c} = log\_b(a)}*

\begin{longtable}[]{@{}
  >{\raggedright\arraybackslash}p{(\columnwidth - 6\tabcolsep) * \real{0.0600}}
  >{\raggedright\arraybackslash}p{(\columnwidth - 6\tabcolsep) * \real{0.4236}}
  >{\raggedright\arraybackslash}p{(\columnwidth - 6\tabcolsep) * \real{0.3385}}
  >{\raggedright\arraybackslash}p{(\columnwidth - 6\tabcolsep) * \real{0.1830}}@{}}
\toprule\noalign{}
\rowcolor{acc!16}
\begin{minipage}[b]{\linewidth}\raggedright
Case
\end{minipage} & \begin{minipage}[b]{\linewidth}\raggedright
Condition
\end{minipage} & \begin{minipage}[b]{\linewidth}\raggedright
Result
\end{minipage} & \begin{minipage}[b]{\linewidth}\raggedright
Intuition
\end{minipage} \\
\midrule\noalign{}
\endhead
\bottomrule\noalign{}
\endlastfoot
Case 1 & f(n) = O(n\^{}(c*\(-\)\(\varepsilon\))) for some \(\varepsilon\) \textgreater{} 0 & T(n) = \(\Theta\)(n\^{}c*) & Leaves dominate \\
\rowcolor{zebra}
Case 2 & f(n) = \(\Theta\)(n\^{}c* \(\cdot\) log\^{}k n), k \(\geq\) 0 & T(n) = \(\Theta\)(n\^{}c* \(\cdot\) log\^{}(k+1) n) & Equal work per level \\
Case 3 & f(n) = \(\Omega\)(n\^{}(c*+\(\varepsilon\))) and regularity holds & T(n) = \(\Theta\)(f(n)) & Root dominates \\
\end{longtable}

\textbf{Master Theorem examples:}

\begin{longtable}[]{@{}
  >{\raggedright\arraybackslash}p{(\columnwidth - 12\tabcolsep) * \real{0.2755}}
  >{\raggedright\arraybackslash}p{(\columnwidth - 12\tabcolsep) * \real{0.0600}}
  >{\raggedright\arraybackslash}p{(\columnwidth - 12\tabcolsep) * \real{0.0600}}
  >{\raggedright\arraybackslash}p{(\columnwidth - 12\tabcolsep) * \real{0.0600}}
  >{\raggedright\arraybackslash}p{(\columnwidth - 12\tabcolsep) * \real{0.2617}}
  >{\raggedright\arraybackslash}p{(\columnwidth - 12\tabcolsep) * \real{0.0634}}
  >{\raggedright\arraybackslash}p{(\columnwidth - 12\tabcolsep) * \real{0.2342}}@{}}
\toprule\noalign{}
\rowcolor{acc!16}
\begin{minipage}[b]{\linewidth}\raggedright
Recurrence
\end{minipage} & \begin{minipage}[b]{\linewidth}\raggedright
a
\end{minipage} & \begin{minipage}[b]{\linewidth}\raggedright
b
\end{minipage} & \begin{minipage}[b]{\linewidth}\raggedright
c*
\end{minipage} & \begin{minipage}[b]{\linewidth}\raggedright
f(n)
\end{minipage} & \begin{minipage}[b]{\linewidth}\raggedright
Case
\end{minipage} & \begin{minipage}[b]{\linewidth}\raggedright
Result
\end{minipage} \\
\midrule\noalign{}
\endhead
\bottomrule\noalign{}
\endlastfoot
T(n) = 2T(n/2) + n & 2 & 2 & 1 & n = \(\Theta\)(n¹) & 2 & \(\Theta\)(n log n) \\
\rowcolor{zebra}
T(n) = 2T(n/2) + 1 & 2 & 2 & 1 & 1 = O(n\^{}0.5) & 1 & \(\Theta\)(n) \\
T(n) = 4T(n/2) + n\(^2\) & 4 & 2 & 2 & n\(^2\) = \(\Theta\)(n\(^2\)) & 2 & \(\Theta\)(n\(^2\) log n) \\
\rowcolor{zebra}
T(n) = 4T(n/2) + n\(^3\) & 4 & 2 & 2 & n\(^3\) = \(\Omega\)(n\^{}2.5) & 3 & \(\Theta\)(n\(^3\)) \\
T(n) = T(n/2) + 1 & 1 & 2 & 0 & 1 = \(\Theta\)(1) & 2 & \(\Theta\)(log n) \\
\rowcolor{zebra}
T(n) = 3T(n/3) + n & 3 & 3 & 1 & n = \(\Theta\)(n¹) & 2 & \(\Theta\)(n log n) \\
\end{longtable}

\textbf{Master Theorem does NOT apply when:}

\begin{itemize}
\tightlist
\item
  Not in form a\(\cdot\)T(n/b) + f(n) --- e.g., T(n) = T(n\(-\)1) + n (linear decrease, not fraction)
\item
  T(n) = T(n/2) + T(n/3) + n (unequal subproblems)
\item
  For T(n) = T(n\(-\)1) + n \(\rightarrow\) sum 1+2+\ldots+n = O(n\(^2\))
\item
  For T(n) = T(n\(-\)1) + 1 \(\rightarrow\) n recursive calls \(\rightarrow\) O(n)
\end{itemize}

\subsection{\texorpdfstring{Input Size \(\rightarrow\) Acceptable Complexity}{Input Size \textbackslash rightarrow Acceptable Complexity}}\label{input-size-rightarrow-acceptable-complexity}

Use this table to instantly identify which algorithms are feasible given constraints.

\begin{longtable}[]{@{}
  >{\raggedright\arraybackslash}p{(\columnwidth - 4\tabcolsep) * \real{0.2054}}
  >{\raggedright\arraybackslash}p{(\columnwidth - 4\tabcolsep) * \real{0.3474}}
  >{\raggedright\arraybackslash}p{(\columnwidth - 4\tabcolsep) * \real{0.4471}}@{}}
\toprule\noalign{}
\rowcolor{acc!16}
\begin{minipage}[b]{\linewidth}\raggedright
n (input size)
\end{minipage} & \begin{minipage}[b]{\linewidth}\raggedright
Max acceptable complexity
\end{minipage} & \begin{minipage}[b]{\linewidth}\raggedright
Typical algorithm
\end{minipage} \\
\midrule\noalign{}
\endhead
\bottomrule\noalign{}
\endlastfoot
n \(\leq\) 10 & O(n!) or O(n\^{}n) & Brute-force, permutations \\
\rowcolor{zebra}
n \(\leq\) 20 & O(2\(^n\)\(\cdot\)n) & Bitmask DP, meet in the middle \\
n \(\leq\) 100 & O(n\(^3\)) or O(n\(^4\)) & Floyd-Warshall, naive string matching \\
\rowcolor{zebra}
n \(\leq\) 1,000 & O(n\(^2\)) & Insertion sort, naive DP \\
n \(\leq\) 10,000 & O(n\(^2\) log n) & Slower O(n\(^2\)) algorithms \\
\rowcolor{zebra}
n \(\leq\) 100,000 & O(n log n) or O(n\(\surd\)n) & Merge sort, segment tree, sqrt decomp \\
n \(\leq\) 1,000,000 & O(n log n) & Merge sort, heap operations \\
\rowcolor{zebra}
n \(\leq\) 10,000,000 & O(n) & Linear scans, counting sort \\
n \textgreater{} 10\^{}8 & O(log n) or O(1) & Binary search, math formulas \\
\end{longtable}

\textbf{Rule of thumb:} Modern computers do \textasciitilde10\^{}8 simple operations/second. Design accordingly.

\section{Algorithms}\label{algorithms}

\subsection{Sorting}\label{sorting}

\subsubsection{Comparison-Based Sorts}\label{comparison-based-sorts}

Comparison sorts have a theoretical lower bound of \textbf{\(\Omega\)(n log n)} --- you cannot do better with only comparison operations.

\textbf{Merge Sort}

\setcodelang{Python}

\begin{Shaded}
\begin{Highlighting}[]
\KeywordTok{def}\NormalTok{ merge\_sort(arr):}
    \ControlFlowTok{if} \BuiltInTok{len}\NormalTok{(arr) }\OperatorTok{\textless{}=} \DecValTok{1}\NormalTok{:}
        \ControlFlowTok{return}\NormalTok{ arr}
\NormalTok{    mid }\OperatorTok{=} \BuiltInTok{len}\NormalTok{(arr) }\OperatorTok{//} \DecValTok{2}
\NormalTok{    left }\OperatorTok{=}\NormalTok{ merge\_sort(arr[:mid])}
\NormalTok{    right }\OperatorTok{=}\NormalTok{ merge\_sort(arr[mid:])}
    \ControlFlowTok{return}\NormalTok{ merge(left, right)}

\KeywordTok{def}\NormalTok{ merge(left, right):}
\NormalTok{    result }\OperatorTok{=}\NormalTok{ []}
\NormalTok{    i }\OperatorTok{=}\NormalTok{ j }\OperatorTok{=} \DecValTok{0}
    \ControlFlowTok{while}\NormalTok{ i }\OperatorTok{\textless{}} \BuiltInTok{len}\NormalTok{(left) }\KeywordTok{and}\NormalTok{ j }\OperatorTok{\textless{}} \BuiltInTok{len}\NormalTok{(right):}
        \ControlFlowTok{if}\NormalTok{ left[i] }\OperatorTok{\textless{}=}\NormalTok{ right[j]:   }\CommentTok{\# \textless{}= preserves stability}
\NormalTok{            result.append(left[i])}\OperatorTok{;}\NormalTok{ i }\OperatorTok{+=} \DecValTok{1}
        \ControlFlowTok{else}\NormalTok{:}
\NormalTok{            result.append(right[j])}\OperatorTok{;}\NormalTok{ j }\OperatorTok{+=} \DecValTok{1}
\NormalTok{    result.extend(left[i:])}
\NormalTok{    result.extend(right[j:])}
    \ControlFlowTok{return}\NormalTok{ result}
\end{Highlighting}
\end{Shaded}

\textbf{Key insight:} Divide conquer and merge. Recurrence T(n) = 2T(n/2) + n \(\rightarrow\) \textbf{O(n log n)} always.
\textbf{Stability:} Yes --- equal elements keep original relative order.
\textbf{Drawback:} O(n) extra space; cache-unfriendly for large arrays.

\textbf{Quick Sort}

\setcodelang{Python}

\begin{Shaded}
\begin{Highlighting}[]
\ImportTok{import}\NormalTok{ random}

\KeywordTok{def}\NormalTok{ quicksort(arr, lo, hi):}
    \ControlFlowTok{if}\NormalTok{ lo }\OperatorTok{\textless{}}\NormalTok{ hi:}
\NormalTok{        p }\OperatorTok{=}\NormalTok{ partition(arr, lo, hi)}
\NormalTok{        quicksort(arr, lo, p }\OperatorTok{{-}} \DecValTok{1}\NormalTok{)}
\NormalTok{        quicksort(arr, p }\OperatorTok{+} \DecValTok{1}\NormalTok{, hi)}

\KeywordTok{def}\NormalTok{ partition(arr, lo, hi):}
    \CommentTok{\# Lomuto partition scheme}
    \CommentTok{\# Randomize pivot to avoid worst{-}case O(n²) on sorted input}
\NormalTok{    rand\_idx }\OperatorTok{=}\NormalTok{ random.randint(lo, hi)}
\NormalTok{    arr[rand\_idx], arr[hi] }\OperatorTok{=}\NormalTok{ arr[hi], arr[rand\_idx]}
\NormalTok{    pivot }\OperatorTok{=}\NormalTok{ arr[hi]}
\NormalTok{    i }\OperatorTok{=}\NormalTok{ lo }\OperatorTok{{-}} \DecValTok{1}
    \ControlFlowTok{for}\NormalTok{ j }\KeywordTok{in} \BuiltInTok{range}\NormalTok{(lo, hi):}
        \ControlFlowTok{if}\NormalTok{ arr[j] }\OperatorTok{\textless{}=}\NormalTok{ pivot:}
\NormalTok{            i }\OperatorTok{+=} \DecValTok{1}
\NormalTok{            arr[i], arr[j] }\OperatorTok{=}\NormalTok{ arr[j], arr[i]}
\NormalTok{    arr[i }\OperatorTok{+} \DecValTok{1}\NormalTok{], arr[hi] }\OperatorTok{=}\NormalTok{ arr[hi], arr[i }\OperatorTok{+} \DecValTok{1}\NormalTok{]}
    \ControlFlowTok{return}\NormalTok{ i }\OperatorTok{+} \DecValTok{1}

\CommentTok{\# Hoare partition (more efficient in practice):}
\KeywordTok{def}\NormalTok{ hoare\_partition(arr, lo, hi):}
\NormalTok{    pivot }\OperatorTok{=}\NormalTok{ arr[(lo }\OperatorTok{+}\NormalTok{ hi) }\OperatorTok{//} \DecValTok{2}\NormalTok{]}
\NormalTok{    i, j }\OperatorTok{=}\NormalTok{ lo }\OperatorTok{{-}} \DecValTok{1}\NormalTok{, hi }\OperatorTok{+} \DecValTok{1}
    \ControlFlowTok{while} \VariableTok{True}\NormalTok{:}
\NormalTok{        i }\OperatorTok{+=} \DecValTok{1}
        \ControlFlowTok{while}\NormalTok{ arr[i] }\OperatorTok{\textless{}}\NormalTok{ pivot: i }\OperatorTok{+=} \DecValTok{1}
\NormalTok{        j }\OperatorTok{{-}=} \DecValTok{1}
        \ControlFlowTok{while}\NormalTok{ arr[j] }\OperatorTok{\textgreater{}}\NormalTok{ pivot: j }\OperatorTok{{-}=} \DecValTok{1}
        \ControlFlowTok{if}\NormalTok{ i }\OperatorTok{\textgreater{}=}\NormalTok{ j: }\ControlFlowTok{return}\NormalTok{ j}
\NormalTok{        arr[i], arr[j] }\OperatorTok{=}\NormalTok{ arr[j], arr[i]}
\end{Highlighting}
\end{Shaded}

\textbf{Key insight:} Pick pivot, partition so left \(\leq\) pivot \(\leq\) right, recurse on both sides.
\textbf{Pivot strategies:}

\begin{itemize}
\tightlist
\item
  Last element: O(n\(^2\)) on sorted/reverse-sorted input
\item
  Random pivot: Expected O(n log n), avoids adversarial inputs
\item
  Median-of-three: Good heuristic
\item
  \textbf{3-way partition (Dutch National Flag):} Best for arrays with many duplicates --- partitions into \texttt{\textless{}\ pivot}, \texttt{==\ pivot}, \texttt{\textgreater{}\ pivot}
\end{itemize}

\textbf{Stability:} No --- partition swaps break relative order.

\textbf{Heap Sort}

\setcodelang{Python}

\begin{Shaded}
\begin{Highlighting}[]
\KeywordTok{def}\NormalTok{ heap\_sort(arr):}
\NormalTok{    n }\OperatorTok{=} \BuiltInTok{len}\NormalTok{(arr)}
    \CommentTok{\# Build max{-}heap (heapify from last internal node)}
    \ControlFlowTok{for}\NormalTok{ i }\KeywordTok{in} \BuiltInTok{range}\NormalTok{(n }\OperatorTok{//} \DecValTok{2} \OperatorTok{{-}} \DecValTok{1}\NormalTok{, }\OperatorTok{{-}}\DecValTok{1}\NormalTok{, }\OperatorTok{{-}}\DecValTok{1}\NormalTok{):}
\NormalTok{        heapify(arr, n, i)}
    \CommentTok{\# Extract elements}
    \ControlFlowTok{for}\NormalTok{ i }\KeywordTok{in} \BuiltInTok{range}\NormalTok{(n }\OperatorTok{{-}} \DecValTok{1}\NormalTok{, }\DecValTok{0}\NormalTok{, }\OperatorTok{{-}}\DecValTok{1}\NormalTok{):}
\NormalTok{        arr[}\DecValTok{0}\NormalTok{], arr[i] }\OperatorTok{=}\NormalTok{ arr[i], arr[}\DecValTok{0}\NormalTok{]  }\CommentTok{\# move max to end}
\NormalTok{        heapify(arr, i, }\DecValTok{0}\NormalTok{)}

\KeywordTok{def}\NormalTok{ heapify(arr, n, i):}
\NormalTok{    largest }\OperatorTok{=}\NormalTok{ i}
\NormalTok{    left, right }\OperatorTok{=} \DecValTok{2}\OperatorTok{*}\NormalTok{i }\OperatorTok{+} \DecValTok{1}\NormalTok{, }\DecValTok{2}\OperatorTok{*}\NormalTok{i }\OperatorTok{+} \DecValTok{2}
    \ControlFlowTok{if}\NormalTok{ left }\OperatorTok{\textless{}}\NormalTok{ n }\KeywordTok{and}\NormalTok{ arr[left] }\OperatorTok{\textgreater{}}\NormalTok{ arr[largest]: largest }\OperatorTok{=}\NormalTok{ left}
    \ControlFlowTok{if}\NormalTok{ right }\OperatorTok{\textless{}}\NormalTok{ n }\KeywordTok{and}\NormalTok{ arr[right] }\OperatorTok{\textgreater{}}\NormalTok{ arr[largest]: largest }\OperatorTok{=}\NormalTok{ right}
    \ControlFlowTok{if}\NormalTok{ largest }\OperatorTok{!=}\NormalTok{ i:}
\NormalTok{        arr[i], arr[largest] }\OperatorTok{=}\NormalTok{ arr[largest], arr[i]}
\NormalTok{        heapify(arr, n, largest)}
\end{Highlighting}
\end{Shaded}

\textbf{Key insight:} Build max-heap in O(n), then repeatedly extract-max in O(log n) each.
\textbf{Stability:} No. \textbf{Space:} O(1) in-place (advantage over merge sort).

\textbf{Insertion Sort}

\setcodelang{Python}

\begin{Shaded}
\begin{Highlighting}[]
\KeywordTok{def}\NormalTok{ insertion\_sort(arr):}
    \ControlFlowTok{for}\NormalTok{ i }\KeywordTok{in} \BuiltInTok{range}\NormalTok{(}\DecValTok{1}\NormalTok{, }\BuiltInTok{len}\NormalTok{(arr)):}
\NormalTok{        key }\OperatorTok{=}\NormalTok{ arr[i]}
\NormalTok{        j }\OperatorTok{=}\NormalTok{ i }\OperatorTok{{-}} \DecValTok{1}
        \ControlFlowTok{while}\NormalTok{ j }\OperatorTok{\textgreater{}=} \DecValTok{0} \KeywordTok{and}\NormalTok{ arr[j] }\OperatorTok{\textgreater{}}\NormalTok{ key:}
\NormalTok{            arr[j }\OperatorTok{+} \DecValTok{1}\NormalTok{] }\OperatorTok{=}\NormalTok{ arr[j]}
\NormalTok{            j }\OperatorTok{{-}=} \DecValTok{1}
\NormalTok{        arr[j }\OperatorTok{+} \DecValTok{1}\NormalTok{] }\OperatorTok{=}\NormalTok{ key}
    \ControlFlowTok{return}\NormalTok{ arr}
\end{Highlighting}
\end{Shaded}

\textbf{Key insight:} Build sorted portion left-to-right by inserting each element in its correct position.
\textbf{Best case:} O(n) when already sorted (inner loop never executes).
\textbf{When to use:} Small arrays (n \textless{} 20), nearly-sorted data, online algorithms (elements arrive one by one).
\textbf{Stability:} Yes.

\subsubsection{Non-Comparison Sorts}\label{non-comparison-sorts}

Non-comparison sorts break the O(n log n) barrier by exploiting structure of keys.

\textbf{Counting Sort}

\setcodelang{Python}

\begin{Shaded}
\begin{Highlighting}[]
\KeywordTok{def}\NormalTok{ counting\_sort(arr, max\_val):}
\NormalTok{    count }\OperatorTok{=}\NormalTok{ [}\DecValTok{0}\NormalTok{] }\OperatorTok{*}\NormalTok{ (max\_val }\OperatorTok{+} \DecValTok{1}\NormalTok{)}
    \ControlFlowTok{for}\NormalTok{ x }\KeywordTok{in}\NormalTok{ arr:}
\NormalTok{        count[x] }\OperatorTok{+=} \DecValTok{1}
    \CommentTok{\# Build cumulative counts for stability}
    \ControlFlowTok{for}\NormalTok{ i }\KeywordTok{in} \BuiltInTok{range}\NormalTok{(}\DecValTok{1}\NormalTok{, max\_val }\OperatorTok{+} \DecValTok{1}\NormalTok{):}
\NormalTok{        count[i] }\OperatorTok{+=}\NormalTok{ count[i }\OperatorTok{{-}} \DecValTok{1}\NormalTok{]}
\NormalTok{    output }\OperatorTok{=}\NormalTok{ [}\DecValTok{0}\NormalTok{] }\OperatorTok{*} \BuiltInTok{len}\NormalTok{(arr)}
    \ControlFlowTok{for}\NormalTok{ x }\KeywordTok{in} \BuiltInTok{reversed}\NormalTok{(arr):     }\CommentTok{\# reverse for stability}
\NormalTok{        output[count[x] }\OperatorTok{{-}} \DecValTok{1}\NormalTok{] }\OperatorTok{=}\NormalTok{ x}
\NormalTok{        count[x] }\OperatorTok{{-}=} \DecValTok{1}
    \ControlFlowTok{return}\NormalTok{ output}
\end{Highlighting}
\end{Shaded}

\textbf{Complexity:} O(n + k) time, O(n + k) space where k = range of values.
\textbf{When to use:} Keys are integers in a known small range. Useless if k \textgreater\textgreater{} n.

\textbf{Radix Sort}

\setcodelang{Python}

\begin{Shaded}
\begin{Highlighting}[]
\KeywordTok{def}\NormalTok{ radix\_sort(arr):}
    \ControlFlowTok{if} \KeywordTok{not}\NormalTok{ arr: }\ControlFlowTok{return}\NormalTok{ arr}
\NormalTok{    max\_val }\OperatorTok{=} \BuiltInTok{max}\NormalTok{(arr)}
\NormalTok{    exp }\OperatorTok{=} \DecValTok{1}
    \ControlFlowTok{while}\NormalTok{ max\_val }\OperatorTok{//}\NormalTok{ exp }\OperatorTok{\textgreater{}} \DecValTok{0}\NormalTok{:}
\NormalTok{        arr }\OperatorTok{=}\NormalTok{ counting\_sort\_by\_digit(arr, exp)}
\NormalTok{        exp }\OperatorTok{*=} \DecValTok{10}
    \ControlFlowTok{return}\NormalTok{ arr}

\KeywordTok{def}\NormalTok{ counting\_sort\_by\_digit(arr, exp):}
\NormalTok{    n }\OperatorTok{=} \BuiltInTok{len}\NormalTok{(arr)}
\NormalTok{    output }\OperatorTok{=}\NormalTok{ [}\DecValTok{0}\NormalTok{] }\OperatorTok{*}\NormalTok{ n}
\NormalTok{    count }\OperatorTok{=}\NormalTok{ [}\DecValTok{0}\NormalTok{] }\OperatorTok{*} \DecValTok{10}
    \ControlFlowTok{for}\NormalTok{ x }\KeywordTok{in}\NormalTok{ arr:}
\NormalTok{        count[(x }\OperatorTok{//}\NormalTok{ exp) }\OperatorTok{\%} \DecValTok{10}\NormalTok{] }\OperatorTok{+=} \DecValTok{1}
    \ControlFlowTok{for}\NormalTok{ i }\KeywordTok{in} \BuiltInTok{range}\NormalTok{(}\DecValTok{1}\NormalTok{, }\DecValTok{10}\NormalTok{):}
\NormalTok{        count[i] }\OperatorTok{+=}\NormalTok{ count[i }\OperatorTok{{-}} \DecValTok{1}\NormalTok{]}
    \ControlFlowTok{for}\NormalTok{ x }\KeywordTok{in} \BuiltInTok{reversed}\NormalTok{(arr):}
\NormalTok{        idx }\OperatorTok{=}\NormalTok{ (x }\OperatorTok{//}\NormalTok{ exp) }\OperatorTok{\%} \DecValTok{10}
\NormalTok{        output[count[idx] }\OperatorTok{{-}} \DecValTok{1}\NormalTok{] }\OperatorTok{=}\NormalTok{ x}
\NormalTok{        count[idx] }\OperatorTok{{-}=} \DecValTok{1}
    \ControlFlowTok{return}\NormalTok{ output}
\end{Highlighting}
\end{Shaded}

\textbf{Complexity:} O(d\(\cdot\)(n + k)) where d = digits, k = base (10 for decimal). Effectively O(n) for fixed-width integers.
\textbf{Stability:} Requires stable inner sort (counting sort). Must process digits from \textbf{least significant} to most significant (LSD radix sort).

\textbf{Bucket Sort}

\setcodelang{Python}

\begin{Shaded}
\begin{Highlighting}[]
\KeywordTok{def}\NormalTok{ bucket\_sort(arr, num\_buckets}\OperatorTok{=}\DecValTok{10}\NormalTok{):}
    \ControlFlowTok{if} \KeywordTok{not}\NormalTok{ arr: }\ControlFlowTok{return}\NormalTok{ arr}
\NormalTok{    min\_val, max\_val }\OperatorTok{=} \BuiltInTok{min}\NormalTok{(arr), }\BuiltInTok{max}\NormalTok{(arr)}
\NormalTok{    bucket\_range }\OperatorTok{=}\NormalTok{ (max\_val }\OperatorTok{{-}}\NormalTok{ min\_val) }\OperatorTok{/}\NormalTok{ num\_buckets }\OperatorTok{+} \FloatTok{1e{-}9}
\NormalTok{    buckets }\OperatorTok{=}\NormalTok{ [[] }\ControlFlowTok{for}\NormalTok{ \_ }\KeywordTok{in} \BuiltInTok{range}\NormalTok{(num\_buckets)]}
    \ControlFlowTok{for}\NormalTok{ x }\KeywordTok{in}\NormalTok{ arr:}
\NormalTok{        idx }\OperatorTok{=} \BuiltInTok{int}\NormalTok{((x }\OperatorTok{{-}}\NormalTok{ min\_val) }\OperatorTok{/}\NormalTok{ bucket\_range)}
\NormalTok{        idx }\OperatorTok{=} \BuiltInTok{min}\NormalTok{(idx, num\_buckets }\OperatorTok{{-}} \DecValTok{1}\NormalTok{)}
\NormalTok{        buckets[idx].append(x)}
\NormalTok{    result }\OperatorTok{=}\NormalTok{ []}
    \ControlFlowTok{for}\NormalTok{ b }\KeywordTok{in}\NormalTok{ buckets:}
\NormalTok{        b.sort()        }\CommentTok{\# insertion sort works well for small buckets}
\NormalTok{        result.extend(b)}
    \ControlFlowTok{return}\NormalTok{ result}
\end{Highlighting}
\end{Shaded}

\textbf{Complexity:} O(n + k) average when data is uniformly distributed; O(n\(^2\)) worst case if all elements land in one bucket.
\textbf{When to use:} Floating-point numbers in a known range, uniformly distributed data.

\subsubsection{Sorting Comparison Table}\label{sorting-comparison-table}

\begin{longtable}[]{@{}
  >{\raggedright\arraybackslash}p{(\columnwidth - 12\tabcolsep) * \real{0.1395}}
  >{\raggedright\arraybackslash}p{(\columnwidth - 12\tabcolsep) * \real{0.0997}}
  >{\raggedright\arraybackslash}p{(\columnwidth - 12\tabcolsep) * \real{0.0997}}
  >{\raggedright\arraybackslash}p{(\columnwidth - 12\tabcolsep) * \real{0.0997}}
  >{\raggedright\arraybackslash}p{(\columnwidth - 12\tabcolsep) * \real{0.0897}}
  >{\raggedright\arraybackslash}p{(\columnwidth - 12\tabcolsep) * \real{0.0688}}
  >{\raggedright\arraybackslash}p{(\columnwidth - 12\tabcolsep) * \real{0.4031}}@{}}
\toprule\noalign{}
\rowcolor{acc!16}
\begin{minipage}[b]{\linewidth}\raggedright
Algorithm
\end{minipage} & \begin{minipage}[b]{\linewidth}\raggedright
Best
\end{minipage} & \begin{minipage}[b]{\linewidth}\raggedright
Average
\end{minipage} & \begin{minipage}[b]{\linewidth}\raggedright
Worst
\end{minipage} & \begin{minipage}[b]{\linewidth}\raggedright
Space
\end{minipage} & \begin{minipage}[b]{\linewidth}\raggedright
Stable
\end{minipage} & \begin{minipage}[b]{\linewidth}\raggedright
Notes
\end{minipage} \\
\midrule\noalign{}
\endhead
\bottomrule\noalign{}
\endlastfoot
Merge Sort & O(n log n) & O(n log n) & O(n log n) & O(n) & Yes & Guaranteed; good for linked lists \\
\rowcolor{zebra}
Quick Sort & O(n log n) & O(n log n) & O(n\(^2\)) & O(log n)* & No & Fastest in practice; randomize pivot \\
Heap Sort & O(n log n) & O(n log n) & O(n log n) & O(1) & No & Cache-unfriendly \\
\rowcolor{zebra}
Insertion Sort & O(n) & O(n\(^2\)) & O(n\(^2\)) & O(1) & Yes & Best for small/nearly-sorted \\
Selection Sort & O(n\(^2\)) & O(n\(^2\)) & O(n\(^2\)) & O(1) & No & Never use in interviews \\
\rowcolor{zebra}
Bubble Sort & O(n) & O(n\(^2\)) & O(n\(^2\)) & O(1) & Yes & Never use in interviews \\
Counting Sort & O(n+k) & O(n+k) & O(n+k) & O(n+k) & Yes & Only integers in small range \\
\rowcolor{zebra}
Radix Sort & O(d\(\cdot\)n) & O(d\(\cdot\)n) & O(d\(\cdot\)n) & O(n+k) & Yes & Fixed-width keys \\
Bucket Sort & O(n+k) & O(n+k) & O(n\(^2\)) & O(n+k) & Yes & Uniform float data \\
\rowcolor{zebra}
Tim Sort & O(n) & O(n log n) & O(n log n) & O(n) & Yes & Python\textquotesingle s built-in; hybrid merge+insertion \\
\end{longtable}

*Quick sort call stack: O(log n) average, O(n) worst.

\textbf{Stability matters when:} You sort by multiple keys sequentially, or the original order of equal elements has meaning (e.g., sort students by grade then by name).

\textbf{When to use which:}

\begin{itemize}
\tightlist
\item
  \textbf{General purpose:} Quick sort (in-place, fast cache) or merge sort (stable)
\item
  \textbf{Guaranteed worst case:} Merge sort or heap sort
\item
  \textbf{Memory-constrained:} Heap sort (O(1) space)
\item
  \textbf{Nearly sorted:} Insertion sort or Tim sort
\item
  \textbf{Integer keys, small range:} Counting sort
\item
  \textbf{Large integers, many of same size:} Radix sort
\item
  \textbf{Floating point, uniform:} Bucket sort
\item
  \textbf{Linked lists:} Merge sort (no random access needed)
\end{itemize}

\subsection{Searching}\label{searching}

\subsubsection{Linear Search}\label{linear-search}

\setcodelang{Python}

\begin{Shaded}
\begin{Highlighting}[]
\KeywordTok{def}\NormalTok{ linear\_search(arr, target):}
    \ControlFlowTok{for}\NormalTok{ i, x }\KeywordTok{in} \BuiltInTok{enumerate}\NormalTok{(arr):}
        \ControlFlowTok{if}\NormalTok{ x }\OperatorTok{==}\NormalTok{ target:}
            \ControlFlowTok{return}\NormalTok{ i}
    \ControlFlowTok{return} \OperatorTok{{-}}\DecValTok{1}
\end{Highlighting}
\end{Shaded}

O(n) time, O(1) space. Use only when array is unsorted or n is tiny.

\subsubsection{Binary Search}\label{binary-search}

\textbf{Standard binary search:}

\setcodelang{Python}

\begin{Shaded}
\begin{Highlighting}[]
\KeywordTok{def}\NormalTok{ binary\_search(arr, target):}
\NormalTok{    lo, hi }\OperatorTok{=} \DecValTok{0}\NormalTok{, }\BuiltInTok{len}\NormalTok{(arr) }\OperatorTok{{-}} \DecValTok{1}
    \ControlFlowTok{while}\NormalTok{ lo }\OperatorTok{\textless{}=}\NormalTok{ hi:}
\NormalTok{        mid }\OperatorTok{=}\NormalTok{ lo }\OperatorTok{+}\NormalTok{ (hi }\OperatorTok{{-}}\NormalTok{ lo) }\OperatorTok{//} \DecValTok{2}   \CommentTok{\# avoids integer overflow}
        \ControlFlowTok{if}\NormalTok{ arr[mid] }\OperatorTok{==}\NormalTok{ target:}
            \ControlFlowTok{return}\NormalTok{ mid}
        \ControlFlowTok{elif}\NormalTok{ arr[mid] }\OperatorTok{\textless{}}\NormalTok{ target:}
\NormalTok{            lo }\OperatorTok{=}\NormalTok{ mid }\OperatorTok{+} \DecValTok{1}
        \ControlFlowTok{else}\NormalTok{:}
\NormalTok{            hi }\OperatorTok{=}\NormalTok{ mid }\OperatorTok{{-}} \DecValTok{1}
    \ControlFlowTok{return} \OperatorTok{{-}}\DecValTok{1}
\end{Highlighting}
\end{Shaded}

\textbf{First occurrence of target:}

\setcodelang{Python}

\begin{Shaded}
\begin{Highlighting}[]
\KeywordTok{def}\NormalTok{ first\_occurrence(arr, target):}
\NormalTok{    lo, hi }\OperatorTok{=} \DecValTok{0}\NormalTok{, }\BuiltInTok{len}\NormalTok{(arr) }\OperatorTok{{-}} \DecValTok{1}
\NormalTok{    result }\OperatorTok{=} \OperatorTok{{-}}\DecValTok{1}
    \ControlFlowTok{while}\NormalTok{ lo }\OperatorTok{\textless{}=}\NormalTok{ hi:}
\NormalTok{        mid }\OperatorTok{=}\NormalTok{ (lo }\OperatorTok{+}\NormalTok{ hi) }\OperatorTok{//} \DecValTok{2}
        \ControlFlowTok{if}\NormalTok{ arr[mid] }\OperatorTok{==}\NormalTok{ target:}
\NormalTok{            result }\OperatorTok{=}\NormalTok{ mid}
\NormalTok{            hi }\OperatorTok{=}\NormalTok{ mid }\OperatorTok{{-}} \DecValTok{1}    \CommentTok{\# keep searching left}
        \ControlFlowTok{elif}\NormalTok{ arr[mid] }\OperatorTok{\textless{}}\NormalTok{ target:}
\NormalTok{            lo }\OperatorTok{=}\NormalTok{ mid }\OperatorTok{+} \DecValTok{1}
        \ControlFlowTok{else}\NormalTok{:}
\NormalTok{            hi }\OperatorTok{=}\NormalTok{ mid }\OperatorTok{{-}} \DecValTok{1}
    \ControlFlowTok{return}\NormalTok{ result}
\end{Highlighting}
\end{Shaded}

\textbf{Last occurrence of target:}

\setcodelang{Python}

\begin{Shaded}
\begin{Highlighting}[]
\KeywordTok{def}\NormalTok{ last\_occurrence(arr, target):}
\NormalTok{    lo, hi }\OperatorTok{=} \DecValTok{0}\NormalTok{, }\BuiltInTok{len}\NormalTok{(arr) }\OperatorTok{{-}} \DecValTok{1}
\NormalTok{    result }\OperatorTok{=} \OperatorTok{{-}}\DecValTok{1}
    \ControlFlowTok{while}\NormalTok{ lo }\OperatorTok{\textless{}=}\NormalTok{ hi:}
\NormalTok{        mid }\OperatorTok{=}\NormalTok{ (lo }\OperatorTok{+}\NormalTok{ hi) }\OperatorTok{//} \DecValTok{2}
        \ControlFlowTok{if}\NormalTok{ arr[mid] }\OperatorTok{==}\NormalTok{ target:}
\NormalTok{            result }\OperatorTok{=}\NormalTok{ mid}
\NormalTok{            lo }\OperatorTok{=}\NormalTok{ mid }\OperatorTok{+} \DecValTok{1}    \CommentTok{\# keep searching right}
        \ControlFlowTok{elif}\NormalTok{ arr[mid] }\OperatorTok{\textless{}}\NormalTok{ target:}
\NormalTok{            lo }\OperatorTok{=}\NormalTok{ mid }\OperatorTok{+} \DecValTok{1}
        \ControlFlowTok{else}\NormalTok{:}
\NormalTok{            hi }\OperatorTok{=}\NormalTok{ mid }\OperatorTok{{-}} \DecValTok{1}
    \ControlFlowTok{return}\NormalTok{ result}
\end{Highlighting}
\end{Shaded}

\textbf{Lower bound (first index where arr{[}i{]} \textgreater= target):}

\setcodelang{Python}

\begin{Shaded}
\begin{Highlighting}[]
\ImportTok{import}\NormalTok{ bisect}

\KeywordTok{def}\NormalTok{ lower\_bound(arr, target):}
    \ControlFlowTok{return}\NormalTok{ bisect.bisect\_left(arr, target)}

\CommentTok{\# Manual implementation:}
\KeywordTok{def}\NormalTok{ lower\_bound\_manual(arr, target):}
\NormalTok{    lo, hi }\OperatorTok{=} \DecValTok{0}\NormalTok{, }\BuiltInTok{len}\NormalTok{(arr)   }\CommentTok{\# hi = len(arr), not len{-}1}
    \ControlFlowTok{while}\NormalTok{ lo }\OperatorTok{\textless{}}\NormalTok{ hi:          }\CommentTok{\# lo \textless{} hi, not lo \textless{}= hi}
\NormalTok{        mid }\OperatorTok{=}\NormalTok{ (lo }\OperatorTok{+}\NormalTok{ hi) }\OperatorTok{//} \DecValTok{2}
        \ControlFlowTok{if}\NormalTok{ arr[mid] }\OperatorTok{\textless{}}\NormalTok{ target:}
\NormalTok{            lo }\OperatorTok{=}\NormalTok{ mid }\OperatorTok{+} \DecValTok{1}
        \ControlFlowTok{else}\NormalTok{:}
\NormalTok{            hi }\OperatorTok{=}\NormalTok{ mid        }\CommentTok{\# not mid{-}1}
    \ControlFlowTok{return}\NormalTok{ lo}
\end{Highlighting}
\end{Shaded}

\textbf{Binary search on rotated sorted array:}

\setcodelang{Python}

\begin{Shaded}
\begin{Highlighting}[]
\KeywordTok{def}\NormalTok{ search\_rotated(arr, target):}
\NormalTok{    lo, hi }\OperatorTok{=} \DecValTok{0}\NormalTok{, }\BuiltInTok{len}\NormalTok{(arr) }\OperatorTok{{-}} \DecValTok{1}
    \ControlFlowTok{while}\NormalTok{ lo }\OperatorTok{\textless{}=}\NormalTok{ hi:}
\NormalTok{        mid }\OperatorTok{=}\NormalTok{ (lo }\OperatorTok{+}\NormalTok{ hi) }\OperatorTok{//} \DecValTok{2}
        \ControlFlowTok{if}\NormalTok{ arr[mid] }\OperatorTok{==}\NormalTok{ target:}
            \ControlFlowTok{return}\NormalTok{ mid}
        \CommentTok{\# Left half is sorted}
        \ControlFlowTok{if}\NormalTok{ arr[lo] }\OperatorTok{\textless{}=}\NormalTok{ arr[mid]:}
            \ControlFlowTok{if}\NormalTok{ arr[lo] }\OperatorTok{\textless{}=}\NormalTok{ target }\OperatorTok{\textless{}}\NormalTok{ arr[mid]:}
\NormalTok{                hi }\OperatorTok{=}\NormalTok{ mid }\OperatorTok{{-}} \DecValTok{1}
            \ControlFlowTok{else}\NormalTok{:}
\NormalTok{                lo }\OperatorTok{=}\NormalTok{ mid }\OperatorTok{+} \DecValTok{1}
        \CommentTok{\# Right half is sorted}
        \ControlFlowTok{else}\NormalTok{:}
            \ControlFlowTok{if}\NormalTok{ arr[mid] }\OperatorTok{\textless{}}\NormalTok{ target }\OperatorTok{\textless{}=}\NormalTok{ arr[hi]:}
\NormalTok{                lo }\OperatorTok{=}\NormalTok{ mid }\OperatorTok{+} \DecValTok{1}
            \ControlFlowTok{else}\NormalTok{:}
\NormalTok{                hi }\OperatorTok{=}\NormalTok{ mid }\OperatorTok{{-}} \DecValTok{1}
    \ControlFlowTok{return} \OperatorTok{{-}}\DecValTok{1}
\end{Highlighting}
\end{Shaded}

\textbf{Binary search on answer space (parametric search):}

The most powerful binary search pattern: instead of searching an array, search for the \textbf{answer itself}.

\setcodelang{Python}

\begin{Shaded}
\begin{Highlighting}[]
\CommentTok{\# Example: Koko eating bananas — find minimum eating speed}
\KeywordTok{def}\NormalTok{ min\_eating\_speed(piles, h):}
    \KeywordTok{def}\NormalTok{ can\_finish(speed):}
        \ControlFlowTok{return} \BuiltInTok{sum}\NormalTok{((p }\OperatorTok{+}\NormalTok{ speed }\OperatorTok{{-}} \DecValTok{1}\NormalTok{) }\OperatorTok{//}\NormalTok{ speed }\ControlFlowTok{for}\NormalTok{ p }\KeywordTok{in}\NormalTok{ piles) }\OperatorTok{\textless{}=}\NormalTok{ h}

\NormalTok{    lo, hi }\OperatorTok{=} \DecValTok{1}\NormalTok{, }\BuiltInTok{max}\NormalTok{(piles)}
    \ControlFlowTok{while}\NormalTok{ lo }\OperatorTok{\textless{}}\NormalTok{ hi:}
\NormalTok{        mid }\OperatorTok{=}\NormalTok{ (lo }\OperatorTok{+}\NormalTok{ hi) }\OperatorTok{//} \DecValTok{2}
        \ControlFlowTok{if}\NormalTok{ can\_finish(mid):}
\NormalTok{            hi }\OperatorTok{=}\NormalTok{ mid        }\CommentTok{\# mid works, try smaller}
        \ControlFlowTok{else}\NormalTok{:}
\NormalTok{            lo }\OperatorTok{=}\NormalTok{ mid }\OperatorTok{+} \DecValTok{1}    \CommentTok{\# too slow, need faster}
    \ControlFlowTok{return}\NormalTok{ lo}

\CommentTok{\# Template: "find minimum X such that condition(X) is True"}
\CommentTok{\# lo = min possible, hi = max possible}
\CommentTok{\# condition is monotone: False...False...True...True}
\end{Highlighting}
\end{Shaded}

\textbf{Binary search complexity:} O(log n) time, O(1) space.

\textbf{Binary search mental model:} Maintain an invariant --- the answer is always in {[}lo, hi{]}. Shrink the range by half each iteration. When lo == hi, that\textquotesingle s the answer.

\subsection{Graph Algorithms}\label{graph-algorithms}

Graphs: G = (V, E) where V = vertices (nodes), E = edges.

\textbf{Representations:}

\begin{itemize}
\tightlist
\item
  \textbf{Adjacency list:} \texttt{graph{[}u{]}\ =\ {[}(v,\ weight),\ ...{]}} --- O(V + E) space, efficient for sparse graphs
\item
  \textbf{Adjacency matrix:} \texttt{graph{[}u{]}{[}v{]}\ =\ weight} --- O(V\(^2\)) space, O(1) edge lookup, efficient for dense graphs
\end{itemize}

\subsubsection{BFS (Breadth-First Search)}\label{bfs-breadth-first-search}

\setcodelang{Python}

\begin{Shaded}
\begin{Highlighting}[]
\ImportTok{from}\NormalTok{ collections }\ImportTok{import}\NormalTok{ deque}

\KeywordTok{def}\NormalTok{ bfs(graph, start):}
\NormalTok{    visited }\OperatorTok{=} \BuiltInTok{set}\NormalTok{([start])}
\NormalTok{    queue }\OperatorTok{=}\NormalTok{ deque([start])}
\NormalTok{    order }\OperatorTok{=}\NormalTok{ []}
    \ControlFlowTok{while}\NormalTok{ queue:}
\NormalTok{        node }\OperatorTok{=}\NormalTok{ queue.popleft()}
\NormalTok{        order.append(node)}
        \ControlFlowTok{for}\NormalTok{ neighbor }\KeywordTok{in}\NormalTok{ graph[node]:}
            \ControlFlowTok{if}\NormalTok{ neighbor }\KeywordTok{not} \KeywordTok{in}\NormalTok{ visited:}
\NormalTok{                visited.add(neighbor)}
\NormalTok{                queue.append(neighbor)}
    \ControlFlowTok{return}\NormalTok{ order}

\CommentTok{\# BFS shortest path (unweighted graph):}
\KeywordTok{def}\NormalTok{ bfs\_shortest\_path(graph, start, end):}
    \ControlFlowTok{if}\NormalTok{ start }\OperatorTok{==}\NormalTok{ end: }\ControlFlowTok{return}\NormalTok{ [start]}
\NormalTok{    visited }\OperatorTok{=}\NormalTok{ \{start: }\VariableTok{None}\NormalTok{\}   }\CommentTok{\# node {-}\textgreater{} parent}
\NormalTok{    queue }\OperatorTok{=}\NormalTok{ deque([start])}
    \ControlFlowTok{while}\NormalTok{ queue:}
\NormalTok{        node }\OperatorTok{=}\NormalTok{ queue.popleft()}
        \ControlFlowTok{for}\NormalTok{ neighbor }\KeywordTok{in}\NormalTok{ graph[node]:}
            \ControlFlowTok{if}\NormalTok{ neighbor }\KeywordTok{not} \KeywordTok{in}\NormalTok{ visited:}
\NormalTok{                visited[neighbor] }\OperatorTok{=}\NormalTok{ node}
                \ControlFlowTok{if}\NormalTok{ neighbor }\OperatorTok{==}\NormalTok{ end:}
\NormalTok{                    path }\OperatorTok{=}\NormalTok{ []}
\NormalTok{                    cur }\OperatorTok{=}\NormalTok{ end}
                    \ControlFlowTok{while}\NormalTok{ cur }\KeywordTok{is} \KeywordTok{not} \VariableTok{None}\NormalTok{:}
\NormalTok{                        path.append(cur)}
\NormalTok{                        cur }\OperatorTok{=}\NormalTok{ visited[cur]}
                    \ControlFlowTok{return}\NormalTok{ path[::}\OperatorTok{{-}}\DecValTok{1}\NormalTok{]}
\NormalTok{                queue.append(neighbor)}
    \ControlFlowTok{return}\NormalTok{ []   }\CommentTok{\# no path}
\end{Highlighting}
\end{Shaded}

\textbf{Complexity:} O(V + E)
\textbf{Use when:} Shortest path in unweighted graph, level-order traversal, connected components, bipartite check.
\textbf{Key property:} Visits nodes in order of distance from source --- guarantees shortest path in unweighted graphs.

\subsubsection{DFS (Depth-First Search)}\label{dfs-depth-first-search}

\setcodelang{Python}

\begin{Shaded}
\begin{Highlighting}[]
\KeywordTok{def}\NormalTok{ dfs\_iterative(graph, start):}
\NormalTok{    visited }\OperatorTok{=} \BuiltInTok{set}\NormalTok{()}
\NormalTok{    stack }\OperatorTok{=}\NormalTok{ [start]}
\NormalTok{    order }\OperatorTok{=}\NormalTok{ []}
    \ControlFlowTok{while}\NormalTok{ stack:}
\NormalTok{        node }\OperatorTok{=}\NormalTok{ stack.pop()}
        \ControlFlowTok{if}\NormalTok{ node }\KeywordTok{in}\NormalTok{ visited:}
            \ControlFlowTok{continue}
\NormalTok{        visited.add(node)}
\NormalTok{        order.append(node)}
        \ControlFlowTok{for}\NormalTok{ neighbor }\KeywordTok{in}\NormalTok{ graph[node]:}
            \ControlFlowTok{if}\NormalTok{ neighbor }\KeywordTok{not} \KeywordTok{in}\NormalTok{ visited:}
\NormalTok{                stack.append(neighbor)}
    \ControlFlowTok{return}\NormalTok{ order}

\KeywordTok{def}\NormalTok{ dfs\_recursive(graph, node, visited}\OperatorTok{=}\VariableTok{None}\NormalTok{):}
    \ControlFlowTok{if}\NormalTok{ visited }\KeywordTok{is} \VariableTok{None}\NormalTok{: visited }\OperatorTok{=} \BuiltInTok{set}\NormalTok{()}
\NormalTok{    visited.add(node)}
    \ControlFlowTok{for}\NormalTok{ neighbor }\KeywordTok{in}\NormalTok{ graph[node]:}
        \ControlFlowTok{if}\NormalTok{ neighbor }\KeywordTok{not} \KeywordTok{in}\NormalTok{ visited:}
\NormalTok{            dfs\_recursive(graph, neighbor, visited)}
    \ControlFlowTok{return}\NormalTok{ visited}
\end{Highlighting}
\end{Shaded}

\textbf{Complexity:} O(V + E)
\textbf{Use when:} Cycle detection, topological sort, connected components, path existence, flood fill.
\textbf{Key property:} Explores as deep as possible before backtracking.

\subsubsection{Topological Sort}\label{topological-sort}

Valid only for \textbf{Directed Acyclic Graphs (DAGs)}.

\textbf{Kahn\textquotesingle s Algorithm (BFS-based):}

\setcodelang{Python}

\begin{Shaded}
\begin{Highlighting}[]
\ImportTok{from}\NormalTok{ collections }\ImportTok{import}\NormalTok{ deque}

\KeywordTok{def}\NormalTok{ topological\_sort\_kahn(graph, n):}
    \CommentTok{"""}
\CommentTok{    graph: adjacency list \{node: [neighbors]\}}
\CommentTok{    n: number of nodes (0 to n{-}1)}
\CommentTok{    Returns topological order, or empty list if cycle exists.}
\CommentTok{    """}
\NormalTok{    in\_degree }\OperatorTok{=}\NormalTok{ [}\DecValTok{0}\NormalTok{] }\OperatorTok{*}\NormalTok{ n}
    \ControlFlowTok{for}\NormalTok{ u }\KeywordTok{in} \BuiltInTok{range}\NormalTok{(n):}
        \ControlFlowTok{for}\NormalTok{ v }\KeywordTok{in}\NormalTok{ graph[u]:}
\NormalTok{            in\_degree[v] }\OperatorTok{+=} \DecValTok{1}

\NormalTok{    queue }\OperatorTok{=}\NormalTok{ deque(u }\ControlFlowTok{for}\NormalTok{ u }\KeywordTok{in} \BuiltInTok{range}\NormalTok{(n) }\ControlFlowTok{if}\NormalTok{ in\_degree[u] }\OperatorTok{==} \DecValTok{0}\NormalTok{)}
\NormalTok{    order }\OperatorTok{=}\NormalTok{ []}

    \ControlFlowTok{while}\NormalTok{ queue:}
\NormalTok{        u }\OperatorTok{=}\NormalTok{ queue.popleft()}
\NormalTok{        order.append(u)}
        \ControlFlowTok{for}\NormalTok{ v }\KeywordTok{in}\NormalTok{ graph[u]:}
\NormalTok{            in\_degree[v] }\OperatorTok{{-}=} \DecValTok{1}
            \ControlFlowTok{if}\NormalTok{ in\_degree[v] }\OperatorTok{==} \DecValTok{0}\NormalTok{:}
\NormalTok{                queue.append(v)}

    \ControlFlowTok{if} \BuiltInTok{len}\NormalTok{(order) }\OperatorTok{!=}\NormalTok{ n:}
        \ControlFlowTok{return}\NormalTok{ []  }\CommentTok{\# cycle detected — not a DAG}
    \ControlFlowTok{return}\NormalTok{ order}
\end{Highlighting}
\end{Shaded}

\textbf{DFS-based topological sort:}

\setcodelang{Python}

\begin{Shaded}
\begin{Highlighting}[]
\KeywordTok{def}\NormalTok{ topological\_sort\_dfs(graph, n):}
\NormalTok{    WHITE, GRAY, BLACK }\OperatorTok{=} \DecValTok{0}\NormalTok{, }\DecValTok{1}\NormalTok{, }\DecValTok{2}
\NormalTok{    color }\OperatorTok{=}\NormalTok{ [WHITE] }\OperatorTok{*}\NormalTok{ n}
\NormalTok{    order }\OperatorTok{=}\NormalTok{ []}
\NormalTok{    has\_cycle }\OperatorTok{=}\NormalTok{ [}\VariableTok{False}\NormalTok{]}

    \KeywordTok{def}\NormalTok{ dfs(u):}
        \ControlFlowTok{if}\NormalTok{ has\_cycle[}\DecValTok{0}\NormalTok{]: }\ControlFlowTok{return}
\NormalTok{        color[u] }\OperatorTok{=}\NormalTok{ GRAY}
        \ControlFlowTok{for}\NormalTok{ v }\KeywordTok{in}\NormalTok{ graph[u]:}
            \ControlFlowTok{if}\NormalTok{ color[v] }\OperatorTok{==}\NormalTok{ GRAY:}
\NormalTok{                has\_cycle[}\DecValTok{0}\NormalTok{] }\OperatorTok{=} \VariableTok{True}
                \ControlFlowTok{return}
            \ControlFlowTok{if}\NormalTok{ color[v] }\OperatorTok{==}\NormalTok{ WHITE:}
\NormalTok{                dfs(v)}
\NormalTok{        color[u] }\OperatorTok{=}\NormalTok{ BLACK}
\NormalTok{        order.append(u)   }\CommentTok{\# add to order AFTER all descendants}

    \ControlFlowTok{for}\NormalTok{ u }\KeywordTok{in} \BuiltInTok{range}\NormalTok{(n):}
        \ControlFlowTok{if}\NormalTok{ color[u] }\OperatorTok{==}\NormalTok{ WHITE:}
\NormalTok{            dfs(u)}

    \ControlFlowTok{if}\NormalTok{ has\_cycle[}\DecValTok{0}\NormalTok{]:}
        \ControlFlowTok{return}\NormalTok{ []}
    \ControlFlowTok{return}\NormalTok{ order[::}\OperatorTok{{-}}\DecValTok{1}\NormalTok{]    }\CommentTok{\# reverse post{-}order = topological order}
\end{Highlighting}
\end{Shaded}

\textbf{Complexity:} O(V + E) for both methods.
\textbf{Kahn vs DFS:}

\begin{itemize}
\tightlist
\item
  Kahn: Intuitive, easy cycle detection (len(order) != n), produces lexicographically smallest order with a min-heap
\item
  DFS: More natural for recursive thinking, easier to add finish-time ordering
\end{itemize}

\textbf{Use cases:} Build systems (Makefile), course prerequisites, task scheduling, dependency resolution.

\subsubsection{Dijkstra\textquotesingle s Algorithm}\label{dijkstras-algorithm}

Shortest path from single source; \textbf{requires non-negative edge weights}.

\setcodelang{Python}

\begin{Shaded}
\begin{Highlighting}[]
\ImportTok{import}\NormalTok{ heapq}

\KeywordTok{def}\NormalTok{ dijkstra(graph, start, n):}
    \CommentTok{"""}
\CommentTok{    graph: adjacency list \{u: [(v, weight), ...]\}}
\CommentTok{    Returns dist[] where dist[v] = shortest distance from start to v.}
\CommentTok{    """}
\NormalTok{    dist }\OperatorTok{=}\NormalTok{ [}\BuiltInTok{float}\NormalTok{(}\StringTok{\textquotesingle{}inf\textquotesingle{}}\NormalTok{)] }\OperatorTok{*}\NormalTok{ n}
\NormalTok{    dist[start] }\OperatorTok{=} \DecValTok{0}
    \CommentTok{\# Min{-}heap: (distance, node)}
\NormalTok{    pq }\OperatorTok{=}\NormalTok{ [(}\DecValTok{0}\NormalTok{, start)]}

    \ControlFlowTok{while}\NormalTok{ pq:}
\NormalTok{        d, u }\OperatorTok{=}\NormalTok{ heapq.heappop(pq)}
        \ControlFlowTok{if}\NormalTok{ d }\OperatorTok{\textgreater{}}\NormalTok{ dist[u]:}
            \ControlFlowTok{continue}    \CommentTok{\# stale entry — already found shorter path}
        \ControlFlowTok{for}\NormalTok{ v, w }\KeywordTok{in}\NormalTok{ graph[u]:}
            \ControlFlowTok{if}\NormalTok{ dist[u] }\OperatorTok{+}\NormalTok{ w }\OperatorTok{\textless{}}\NormalTok{ dist[v]:}
\NormalTok{                dist[v] }\OperatorTok{=}\NormalTok{ dist[u] }\OperatorTok{+}\NormalTok{ w}
\NormalTok{                heapq.heappush(pq, (dist[v], v))}

    \ControlFlowTok{return}\NormalTok{ dist}

\CommentTok{\# With path reconstruction:}
\KeywordTok{def}\NormalTok{ dijkstra\_with\_path(graph, start, end, n):}
\NormalTok{    dist }\OperatorTok{=}\NormalTok{ [}\BuiltInTok{float}\NormalTok{(}\StringTok{\textquotesingle{}inf\textquotesingle{}}\NormalTok{)] }\OperatorTok{*}\NormalTok{ n}
\NormalTok{    prev }\OperatorTok{=}\NormalTok{ [}\OperatorTok{{-}}\DecValTok{1}\NormalTok{] }\OperatorTok{*}\NormalTok{ n}
\NormalTok{    dist[start] }\OperatorTok{=} \DecValTok{0}
\NormalTok{    pq }\OperatorTok{=}\NormalTok{ [(}\DecValTok{0}\NormalTok{, start)]}
    \ControlFlowTok{while}\NormalTok{ pq:}
\NormalTok{        d, u }\OperatorTok{=}\NormalTok{ heapq.heappop(pq)}
        \ControlFlowTok{if}\NormalTok{ d }\OperatorTok{\textgreater{}}\NormalTok{ dist[u]: }\ControlFlowTok{continue}
        \ControlFlowTok{if}\NormalTok{ u }\OperatorTok{==}\NormalTok{ end: }\ControlFlowTok{break}
        \ControlFlowTok{for}\NormalTok{ v, w }\KeywordTok{in}\NormalTok{ graph[u]:}
            \ControlFlowTok{if}\NormalTok{ dist[u] }\OperatorTok{+}\NormalTok{ w }\OperatorTok{\textless{}}\NormalTok{ dist[v]:}
\NormalTok{                dist[v] }\OperatorTok{=}\NormalTok{ dist[u] }\OperatorTok{+}\NormalTok{ w}
\NormalTok{                prev[v] }\OperatorTok{=}\NormalTok{ u}
\NormalTok{                heapq.heappush(pq, (dist[v], v))}
    \CommentTok{\# Reconstruct path}
\NormalTok{    path }\OperatorTok{=}\NormalTok{ []}
\NormalTok{    cur }\OperatorTok{=}\NormalTok{ end}
    \ControlFlowTok{while}\NormalTok{ cur }\OperatorTok{!=} \OperatorTok{{-}}\DecValTok{1}\NormalTok{:}
\NormalTok{        path.append(cur)}
\NormalTok{        cur }\OperatorTok{=}\NormalTok{ prev[cur]}
    \ControlFlowTok{return}\NormalTok{ dist[end], path[::}\OperatorTok{{-}}\DecValTok{1}\NormalTok{]}
\end{Highlighting}
\end{Shaded}

\textbf{Complexity:} O((V + E) log V) with binary heap, O(V\(^2\) + E) with array (dense graphs).
\textbf{When to use:} Single-source shortest path, non-negative weights.
\textbf{Does NOT work with negative edges} --- use Bellman-Ford instead.

\subsubsection{Bellman-Ford Algorithm}\label{bellman-ford-algorithm}

Shortest path from single source; \textbf{handles negative edge weights}; detects negative cycles.

\setcodelang{Python}

\begin{Shaded}
\begin{Highlighting}[]
\KeywordTok{def}\NormalTok{ bellman\_ford(edges, n, start):}
    \CommentTok{"""}
\CommentTok{    edges: list of (u, v, weight)}
\CommentTok{    Returns dist[] or None if negative cycle reachable from start.}
\CommentTok{    """}
\NormalTok{    dist }\OperatorTok{=}\NormalTok{ [}\BuiltInTok{float}\NormalTok{(}\StringTok{\textquotesingle{}inf\textquotesingle{}}\NormalTok{)] }\OperatorTok{*}\NormalTok{ n}
\NormalTok{    dist[start] }\OperatorTok{=} \DecValTok{0}

    \CommentTok{\# Relax all edges V{-}1 times}
    \ControlFlowTok{for}\NormalTok{ \_ }\KeywordTok{in} \BuiltInTok{range}\NormalTok{(n }\OperatorTok{{-}} \DecValTok{1}\NormalTok{):}
\NormalTok{        updated }\OperatorTok{=} \VariableTok{False}
        \ControlFlowTok{for}\NormalTok{ u, v, w }\KeywordTok{in}\NormalTok{ edges:}
            \ControlFlowTok{if}\NormalTok{ dist[u] }\OperatorTok{!=} \BuiltInTok{float}\NormalTok{(}\StringTok{\textquotesingle{}inf\textquotesingle{}}\NormalTok{) }\KeywordTok{and}\NormalTok{ dist[u] }\OperatorTok{+}\NormalTok{ w }\OperatorTok{\textless{}}\NormalTok{ dist[v]:}
\NormalTok{                dist[v] }\OperatorTok{=}\NormalTok{ dist[u] }\OperatorTok{+}\NormalTok{ w}
\NormalTok{                updated }\OperatorTok{=} \VariableTok{True}
        \ControlFlowTok{if} \KeywordTok{not}\NormalTok{ updated:}
            \ControlFlowTok{break}   \CommentTok{\# converged early}

    \CommentTok{\# Check for negative cycles (V{-}th relaxation)}
    \ControlFlowTok{for}\NormalTok{ u, v, w }\KeywordTok{in}\NormalTok{ edges:}
        \ControlFlowTok{if}\NormalTok{ dist[u] }\OperatorTok{!=} \BuiltInTok{float}\NormalTok{(}\StringTok{\textquotesingle{}inf\textquotesingle{}}\NormalTok{) }\KeywordTok{and}\NormalTok{ dist[u] }\OperatorTok{+}\NormalTok{ w }\OperatorTok{\textless{}}\NormalTok{ dist[v]:}
            \ControlFlowTok{return} \VariableTok{None}  \CommentTok{\# negative cycle detected}

    \ControlFlowTok{return}\NormalTok{ dist}
\end{Highlighting}
\end{Shaded}

\textbf{Complexity:} O(V\(\cdot\)E)
\textbf{When to use:} Negative edge weights, detecting negative cycles (e.g., currency arbitrage).
\textbf{Guaranteed correct after V\(-\)1 iterations} because any shortest path visits at most V\(-\)1 edges (in a graph with no negative cycles).

\subsubsection{Floyd-Warshall Algorithm}\label{floyd-warshall-algorithm}

All-pairs shortest paths; handles negative edges; detects negative cycles.

\setcodelang{Python}

\begin{Shaded}
\begin{Highlighting}[]
\KeywordTok{def}\NormalTok{ floyd\_warshall(n, edges):}
    \CommentTok{"""Returns dist[i][j] = shortest path from i to j."""}
\NormalTok{    INF }\OperatorTok{=} \BuiltInTok{float}\NormalTok{(}\StringTok{\textquotesingle{}inf\textquotesingle{}}\NormalTok{)}
\NormalTok{    dist }\OperatorTok{=}\NormalTok{ [[INF] }\OperatorTok{*}\NormalTok{ n }\ControlFlowTok{for}\NormalTok{ \_ }\KeywordTok{in} \BuiltInTok{range}\NormalTok{(n)]}
    \ControlFlowTok{for}\NormalTok{ i }\KeywordTok{in} \BuiltInTok{range}\NormalTok{(n):}
\NormalTok{        dist[i][i] }\OperatorTok{=} \DecValTok{0}
    \ControlFlowTok{for}\NormalTok{ u, v, w }\KeywordTok{in}\NormalTok{ edges:}
\NormalTok{        dist[u][v] }\OperatorTok{=} \BuiltInTok{min}\NormalTok{(dist[u][v], w)   }\CommentTok{\# handle multi{-}edges}

    \ControlFlowTok{for}\NormalTok{ k }\KeywordTok{in} \BuiltInTok{range}\NormalTok{(n):          }\CommentTok{\# intermediate vertex}
        \ControlFlowTok{for}\NormalTok{ i }\KeywordTok{in} \BuiltInTok{range}\NormalTok{(n):}
            \ControlFlowTok{for}\NormalTok{ j }\KeywordTok{in} \BuiltInTok{range}\NormalTok{(n):}
                \ControlFlowTok{if}\NormalTok{ dist[i][k] }\OperatorTok{+}\NormalTok{ dist[k][j] }\OperatorTok{\textless{}}\NormalTok{ dist[i][j]:}
\NormalTok{                    dist[i][j] }\OperatorTok{=}\NormalTok{ dist[i][k] }\OperatorTok{+}\NormalTok{ dist[k][j]}

    \CommentTok{\# Check negative cycles: dist[i][i] \textless{} 0 for some i}
    \ControlFlowTok{for}\NormalTok{ i }\KeywordTok{in} \BuiltInTok{range}\NormalTok{(n):}
        \ControlFlowTok{if}\NormalTok{ dist[i][i] }\OperatorTok{\textless{}} \DecValTok{0}\NormalTok{:}
            \ControlFlowTok{return} \VariableTok{None}  \CommentTok{\# negative cycle}

    \ControlFlowTok{return}\NormalTok{ dist}
\end{Highlighting}
\end{Shaded}

\textbf{Complexity:} O(V\(^3\)) time, O(V\(^2\)) space.
\textbf{When to use:} All-pairs shortest paths, transitive closure, detecting negative cycles.

\subsubsection{Graph Algorithm Comparison Table}\label{graph-algorithm-comparison-table}

\begin{longtable}[]{@{}
  >{\raggedright\arraybackslash}p{(\columnwidth - 10\tabcolsep) * \real{0.1171}}
  >{\raggedright\arraybackslash}p{(\columnwidth - 10\tabcolsep) * \real{0.1742}}
  >{\raggedright\arraybackslash}p{(\columnwidth - 10\tabcolsep) * \real{0.1822}}
  >{\raggedright\arraybackslash}p{(\columnwidth - 10\tabcolsep) * \real{0.1267}}
  >{\raggedright\arraybackslash}p{(\columnwidth - 10\tabcolsep) * \real{0.0964}}
  >{\raggedright\arraybackslash}p{(\columnwidth - 10\tabcolsep) * \real{0.3034}}@{}}
\toprule\noalign{}
\rowcolor{acc!16}
\begin{minipage}[b]{\linewidth}\raggedright
Algorithm
\end{minipage} & \begin{minipage}[b]{\linewidth}\raggedright
Handles Negative Edges
\end{minipage} & \begin{minipage}[b]{\linewidth}\raggedright
Handles Negative Cycles
\end{minipage} & \begin{minipage}[b]{\linewidth}\raggedright
Single/All Pairs
\end{minipage} & \begin{minipage}[b]{\linewidth}\raggedright
Complexity
\end{minipage} & \begin{minipage}[b]{\linewidth}\raggedright
Notes
\end{minipage} \\
\midrule\noalign{}
\endhead
\bottomrule\noalign{}
\endlastfoot
Dijkstra & No & No & Single source & O((V+E) log V) & Fastest for non-negative; use with priority queue \\
\rowcolor{zebra}
Bellman-Ford & Yes & Detects & Single source & O(V\(\cdot\)E) & Use when negative edges exist \\
Floyd-Warshall & Yes & Detects & All pairs & O(V\(^3\)) & Best for dense graphs, all-pairs \\
\rowcolor{zebra}
BFS & No (unit weights only) & No & Single source & O(V+E) & Optimal for unweighted graphs \\
DAG Shortest Path & Yes (DAG only) & N/A & Single source & O(V+E) & Topo sort + relax; fastest for DAGs \\
\end{longtable}

\subsubsection{Minimum Spanning Tree (MST)}\label{minimum-spanning-tree-mst}

MST of a connected weighted undirected graph: a spanning tree with minimum total edge weight.

\textbf{Kruskal\textquotesingle s Algorithm:}

\setcodelang{Python}

\begin{Shaded}
\begin{Highlighting}[]
\KeywordTok{def}\NormalTok{ kruskal(n, edges):}
    \CommentTok{"""}
\CommentTok{    edges: list of (weight, u, v)}
\CommentTok{    Returns MST weight and edges.}
\CommentTok{    """}
\NormalTok{    edges.sort()   }\CommentTok{\# sort by weight}
\NormalTok{    parent }\OperatorTok{=} \BuiltInTok{list}\NormalTok{(}\BuiltInTok{range}\NormalTok{(n))}
\NormalTok{    rank }\OperatorTok{=}\NormalTok{ [}\DecValTok{0}\NormalTok{] }\OperatorTok{*}\NormalTok{ n}

    \KeywordTok{def}\NormalTok{ find(x):}
        \ControlFlowTok{while}\NormalTok{ parent[x] }\OperatorTok{!=}\NormalTok{ x:}
\NormalTok{            parent[x] }\OperatorTok{=}\NormalTok{ parent[parent[x]]  }\CommentTok{\# path compression (2{-}step)}
\NormalTok{            x }\OperatorTok{=}\NormalTok{ parent[x]}
        \ControlFlowTok{return}\NormalTok{ x}

    \KeywordTok{def}\NormalTok{ union(x, y):}
\NormalTok{        px, py }\OperatorTok{=}\NormalTok{ find(x), find(y)}
        \ControlFlowTok{if}\NormalTok{ px }\OperatorTok{==}\NormalTok{ py: }\ControlFlowTok{return} \VariableTok{False}   \CommentTok{\# same component — cycle}
        \ControlFlowTok{if}\NormalTok{ rank[px] }\OperatorTok{\textless{}}\NormalTok{ rank[py]:}
\NormalTok{            px, py }\OperatorTok{=}\NormalTok{ py, px}
\NormalTok{        parent[py] }\OperatorTok{=}\NormalTok{ px}
        \ControlFlowTok{if}\NormalTok{ rank[px] }\OperatorTok{==}\NormalTok{ rank[py]:}
\NormalTok{            rank[px] }\OperatorTok{+=} \DecValTok{1}
        \ControlFlowTok{return} \VariableTok{True}

\NormalTok{    mst\_weight }\OperatorTok{=} \DecValTok{0}
\NormalTok{    mst\_edges }\OperatorTok{=}\NormalTok{ []}
    \ControlFlowTok{for}\NormalTok{ w, u, v }\KeywordTok{in}\NormalTok{ edges:}
        \ControlFlowTok{if}\NormalTok{ union(u, v):}
\NormalTok{            mst\_weight }\OperatorTok{+=}\NormalTok{ w}
\NormalTok{            mst\_edges.append((u, v, w))}
            \ControlFlowTok{if} \BuiltInTok{len}\NormalTok{(mst\_edges) }\OperatorTok{==}\NormalTok{ n }\OperatorTok{{-}} \DecValTok{1}\NormalTok{:}
                \ControlFlowTok{break}

    \ControlFlowTok{return}\NormalTok{ mst\_weight, mst\_edges}
\end{Highlighting}
\end{Shaded}

\textbf{Prim\textquotesingle s Algorithm:}

\setcodelang{Python}

\begin{Shaded}
\begin{Highlighting}[]
\ImportTok{import}\NormalTok{ heapq}

\KeywordTok{def}\NormalTok{ prim(graph, n):}
    \CommentTok{"""}
\CommentTok{    graph: adjacency list \{u: [(v, weight), ...]\}}
\CommentTok{    Returns MST weight.}
\CommentTok{    """}
\NormalTok{    visited }\OperatorTok{=} \BuiltInTok{set}\NormalTok{()}
    \CommentTok{\# Start from node 0: (weight, node)}
\NormalTok{    pq }\OperatorTok{=}\NormalTok{ [(}\DecValTok{0}\NormalTok{, }\DecValTok{0}\NormalTok{)]}
\NormalTok{    mst\_weight }\OperatorTok{=} \DecValTok{0}

    \ControlFlowTok{while}\NormalTok{ pq }\KeywordTok{and} \BuiltInTok{len}\NormalTok{(visited) }\OperatorTok{\textless{}}\NormalTok{ n:}
\NormalTok{        w, u }\OperatorTok{=}\NormalTok{ heapq.heappop(pq)}
        \ControlFlowTok{if}\NormalTok{ u }\KeywordTok{in}\NormalTok{ visited:}
            \ControlFlowTok{continue}
\NormalTok{        visited.add(u)}
\NormalTok{        mst\_weight }\OperatorTok{+=}\NormalTok{ w}
        \ControlFlowTok{for}\NormalTok{ v, weight }\KeywordTok{in}\NormalTok{ graph[u]:}
            \ControlFlowTok{if}\NormalTok{ v }\KeywordTok{not} \KeywordTok{in}\NormalTok{ visited:}
\NormalTok{                heapq.heappush(pq, (weight, v))}

    \ControlFlowTok{return}\NormalTok{ mst\_weight }\ControlFlowTok{if} \BuiltInTok{len}\NormalTok{(visited) }\OperatorTok{==}\NormalTok{ n }\ControlFlowTok{else} \OperatorTok{{-}}\DecValTok{1}  \CommentTok{\# {-}1 if disconnected}
\end{Highlighting}
\end{Shaded}

\textbf{Kruskal vs Prim:}
\textbar{} \textbar{} Kruskal \textbar{} Prim \textbar{}
\textbar-\textbar-\/-\/-\/-\/-\/-\/-\/-\/-\textbar-\/-\/-\/-\/-\/-\textbar{}
\textbar{} Approach \textbar{} Edge-centric: sort all edges, add if no cycle \textbar{} Vertex-centric: grow tree from one vertex \textbar{}
\textbar{} Complexity \textbar{} O(E log E) \textbar{} O((V+E) log V) with heap \textbar{}
\textbar{} Best for \textbar{} Sparse graphs \textbar{} Dense graphs \textbar{}
\textbar{} Data structure \textbar{} Union-Find \textbar{} Priority queue \textbar{}

\subsubsection{Union-Find (Disjoint Set Union)}\label{union-find-disjoint-set-union}

\setcodelang{Python}

\begin{Shaded}
\begin{Highlighting}[]
\KeywordTok{class}\NormalTok{ UnionFind:}
    \KeywordTok{def} \FunctionTok{\_\_init\_\_}\NormalTok{(}\VariableTok{self}\NormalTok{, n):}
        \VariableTok{self}\NormalTok{.parent }\OperatorTok{=} \BuiltInTok{list}\NormalTok{(}\BuiltInTok{range}\NormalTok{(n))}
        \VariableTok{self}\NormalTok{.rank }\OperatorTok{=}\NormalTok{ [}\DecValTok{0}\NormalTok{] }\OperatorTok{*}\NormalTok{ n}
        \VariableTok{self}\NormalTok{.components }\OperatorTok{=}\NormalTok{ n}

    \KeywordTok{def}\NormalTok{ find(}\VariableTok{self}\NormalTok{, x):}
        \ControlFlowTok{if} \VariableTok{self}\NormalTok{.parent[x] }\OperatorTok{!=}\NormalTok{ x:}
            \VariableTok{self}\NormalTok{.parent[x] }\OperatorTok{=} \VariableTok{self}\NormalTok{.find(}\VariableTok{self}\NormalTok{.parent[x])   }\CommentTok{\# path compression}
        \ControlFlowTok{return} \VariableTok{self}\NormalTok{.parent[x]}

    \KeywordTok{def}\NormalTok{ union(}\VariableTok{self}\NormalTok{, x, y):}
\NormalTok{        px, py }\OperatorTok{=} \VariableTok{self}\NormalTok{.find(x), }\VariableTok{self}\NormalTok{.find(y)}
        \ControlFlowTok{if}\NormalTok{ px }\OperatorTok{==}\NormalTok{ py: }\ControlFlowTok{return} \VariableTok{False}
        \CommentTok{\# Union by rank}
        \ControlFlowTok{if} \VariableTok{self}\NormalTok{.rank[px] }\OperatorTok{\textless{}} \VariableTok{self}\NormalTok{.rank[py]:}
\NormalTok{            px, py }\OperatorTok{=}\NormalTok{ py, px}
        \VariableTok{self}\NormalTok{.parent[py] }\OperatorTok{=}\NormalTok{ px}
        \ControlFlowTok{if} \VariableTok{self}\NormalTok{.rank[px] }\OperatorTok{==} \VariableTok{self}\NormalTok{.rank[py]:}
            \VariableTok{self}\NormalTok{.rank[px] }\OperatorTok{+=} \DecValTok{1}
        \VariableTok{self}\NormalTok{.components }\OperatorTok{{-}=} \DecValTok{1}
        \ControlFlowTok{return} \VariableTok{True}

    \KeywordTok{def}\NormalTok{ connected(}\VariableTok{self}\NormalTok{, x, y):}
        \ControlFlowTok{return} \VariableTok{self}\NormalTok{.find(x) }\OperatorTok{==} \VariableTok{self}\NormalTok{.find(y)}
\end{Highlighting}
\end{Shaded}

\textbf{Amortized complexity:} O(\(\alpha\)(n)) per operation where \(\alpha\) is inverse Ackermann --- effectively O(1).
\textbf{Use for:} Cycle detection, Kruskal\textquotesingle s MST, connected components, dynamic connectivity.

\subsubsection{Cycle Detection}\label{cycle-detection}

\textbf{Undirected graph --- Union-Find:}

\setcodelang{Python}

\begin{Shaded}
\begin{Highlighting}[]
\KeywordTok{def}\NormalTok{ has\_cycle\_undirected(n, edges):}
\NormalTok{    uf }\OperatorTok{=}\NormalTok{ UnionFind(n)}
    \ControlFlowTok{for}\NormalTok{ u, v }\KeywordTok{in}\NormalTok{ edges:}
        \ControlFlowTok{if} \KeywordTok{not}\NormalTok{ uf.union(u, v):}
            \ControlFlowTok{return} \VariableTok{True}   \CommentTok{\# u and v already in same component}
    \ControlFlowTok{return} \VariableTok{False}
\end{Highlighting}
\end{Shaded}

\textbf{Undirected graph --- DFS (coloring):}

\setcodelang{Python}

\begin{Shaded}
\begin{Highlighting}[]
\KeywordTok{def}\NormalTok{ has\_cycle\_undirected\_dfs(graph, n):}
\NormalTok{    visited }\OperatorTok{=} \BuiltInTok{set}\NormalTok{()}
    \KeywordTok{def}\NormalTok{ dfs(node, parent):}
\NormalTok{        visited.add(node)}
        \ControlFlowTok{for}\NormalTok{ neighbor }\KeywordTok{in}\NormalTok{ graph[node]:}
            \ControlFlowTok{if}\NormalTok{ neighbor }\KeywordTok{not} \KeywordTok{in}\NormalTok{ visited:}
                \ControlFlowTok{if}\NormalTok{ dfs(neighbor, node): }\ControlFlowTok{return} \VariableTok{True}
            \ControlFlowTok{elif}\NormalTok{ neighbor }\OperatorTok{!=}\NormalTok{ parent:    }\CommentTok{\# back edge to non{-}parent = cycle}
                \ControlFlowTok{return} \VariableTok{True}
        \ControlFlowTok{return} \VariableTok{False}
    \ControlFlowTok{for}\NormalTok{ i }\KeywordTok{in} \BuiltInTok{range}\NormalTok{(n):}
        \ControlFlowTok{if}\NormalTok{ i }\KeywordTok{not} \KeywordTok{in}\NormalTok{ visited:}
            \ControlFlowTok{if}\NormalTok{ dfs(i, }\OperatorTok{{-}}\DecValTok{1}\NormalTok{): }\ControlFlowTok{return} \VariableTok{True}
    \ControlFlowTok{return} \VariableTok{False}
\end{Highlighting}
\end{Shaded}

\textbf{Directed graph --- DFS with 3-color:}

\setcodelang{Python}

\begin{Shaded}
\begin{Highlighting}[]
\KeywordTok{def}\NormalTok{ has\_cycle\_directed(graph, n):}
\NormalTok{    WHITE, GRAY, BLACK }\OperatorTok{=} \DecValTok{0}\NormalTok{, }\DecValTok{1}\NormalTok{, }\DecValTok{2}
\NormalTok{    color }\OperatorTok{=}\NormalTok{ [WHITE] }\OperatorTok{*}\NormalTok{ n}
    \KeywordTok{def}\NormalTok{ dfs(u):}
\NormalTok{        color[u] }\OperatorTok{=}\NormalTok{ GRAY}
        \ControlFlowTok{for}\NormalTok{ v }\KeywordTok{in}\NormalTok{ graph[u]:}
            \ControlFlowTok{if}\NormalTok{ color[v] }\OperatorTok{==}\NormalTok{ GRAY: }\ControlFlowTok{return} \VariableTok{True}    \CommentTok{\# back edge = cycle}
            \ControlFlowTok{if}\NormalTok{ color[v] }\OperatorTok{==}\NormalTok{ WHITE }\KeywordTok{and}\NormalTok{ dfs(v): }\ControlFlowTok{return} \VariableTok{True}
\NormalTok{        color[u] }\OperatorTok{=}\NormalTok{ BLACK}
        \ControlFlowTok{return} \VariableTok{False}
    \ControlFlowTok{return} \BuiltInTok{any}\NormalTok{(dfs(u) }\ControlFlowTok{for}\NormalTok{ u }\KeywordTok{in} \BuiltInTok{range}\NormalTok{(n) }\ControlFlowTok{if}\NormalTok{ color[u] }\OperatorTok{==}\NormalTok{ WHITE)}
\end{Highlighting}
\end{Shaded}

\textbf{Key difference:} In undirected graphs, visiting a neighbor (other than parent) that\textquotesingle s already visited = cycle. In directed graphs, visiting a GRAY (currently-being-processed) node = cycle (back edge).

\subsubsection{Connected Components}\label{connected-components}

\setcodelang{Python}

\begin{Shaded}
\begin{Highlighting}[]
\KeywordTok{def}\NormalTok{ connected\_components(graph, n):}
\NormalTok{    visited }\OperatorTok{=} \BuiltInTok{set}\NormalTok{()}
\NormalTok{    components }\OperatorTok{=}\NormalTok{ []}
    \KeywordTok{def}\NormalTok{ dfs(node, component):}
\NormalTok{        visited.add(node)}
\NormalTok{        component.append(node)}
        \ControlFlowTok{for}\NormalTok{ neighbor }\KeywordTok{in}\NormalTok{ graph[node]:}
            \ControlFlowTok{if}\NormalTok{ neighbor }\KeywordTok{not} \KeywordTok{in}\NormalTok{ visited:}
\NormalTok{                dfs(neighbor, component)}
    \ControlFlowTok{for}\NormalTok{ i }\KeywordTok{in} \BuiltInTok{range}\NormalTok{(n):}
        \ControlFlowTok{if}\NormalTok{ i }\KeywordTok{not} \KeywordTok{in}\NormalTok{ visited:}
\NormalTok{            component }\OperatorTok{=}\NormalTok{ []}
\NormalTok{            dfs(i, component)}
\NormalTok{            components.append(component)}
    \ControlFlowTok{return}\NormalTok{ components}
\end{Highlighting}
\end{Shaded}

\subsubsection{Bipartite Check}\label{bipartite-check}

\setcodelang{Python}

\begin{Shaded}
\begin{Highlighting}[]
\ImportTok{from}\NormalTok{ collections }\ImportTok{import}\NormalTok{ deque}

\KeywordTok{def}\NormalTok{ is\_bipartite(graph, n):}
\NormalTok{    color }\OperatorTok{=}\NormalTok{ [}\OperatorTok{{-}}\DecValTok{1}\NormalTok{] }\OperatorTok{*}\NormalTok{ n}
    \ControlFlowTok{for}\NormalTok{ start }\KeywordTok{in} \BuiltInTok{range}\NormalTok{(n):}
        \ControlFlowTok{if}\NormalTok{ color[start] }\OperatorTok{!=} \OperatorTok{{-}}\DecValTok{1}\NormalTok{: }\ControlFlowTok{continue}
\NormalTok{        color[start] }\OperatorTok{=} \DecValTok{0}
\NormalTok{        queue }\OperatorTok{=}\NormalTok{ deque([start])}
        \ControlFlowTok{while}\NormalTok{ queue:}
\NormalTok{            u }\OperatorTok{=}\NormalTok{ queue.popleft()}
            \ControlFlowTok{for}\NormalTok{ v }\KeywordTok{in}\NormalTok{ graph[u]:}
                \ControlFlowTok{if}\NormalTok{ color[v] }\OperatorTok{==} \OperatorTok{{-}}\DecValTok{1}\NormalTok{:}
\NormalTok{                    color[v] }\OperatorTok{=} \DecValTok{1} \OperatorTok{{-}}\NormalTok{ color[u]}
\NormalTok{                    queue.append(v)}
                \ControlFlowTok{elif}\NormalTok{ color[v] }\OperatorTok{==}\NormalTok{ color[u]:}
                    \ControlFlowTok{return} \VariableTok{False}    \CommentTok{\# same color = odd cycle = not bipartite}
    \ControlFlowTok{return} \VariableTok{True}
\end{Highlighting}
\end{Shaded}

\textbf{A graph is bipartite iff it has no odd-length cycles.} Two-color with BFS; if conflict found, not bipartite.

\subsection{Dynamic Programming}\label{dynamic-programming}

DP solves problems by \textbf{breaking them into overlapping subproblems} and storing results to avoid recomputation.

\subsubsection{Memoization vs Tabulation}\label{memoization-vs-tabulation}

\begin{longtable}[]{@{}
  >{\raggedright\arraybackslash}p{(\columnwidth - 4\tabcolsep) * \real{0.1528}}
  >{\raggedright\arraybackslash}p{(\columnwidth - 4\tabcolsep) * \real{0.4306}}
  >{\raggedright\arraybackslash}p{(\columnwidth - 4\tabcolsep) * \real{0.4167}}@{}}
\toprule\noalign{}
\rowcolor{acc!16}
\begin{minipage}[b]{\linewidth}\raggedright
\end{minipage} & \begin{minipage}[b]{\linewidth}\raggedright
Memoization (Top-Down)
\end{minipage} & \begin{minipage}[b]{\linewidth}\raggedright
Tabulation (Bottom-Up)
\end{minipage} \\
\midrule\noalign{}
\endhead
\bottomrule\noalign{}
\endlastfoot
Approach & Recursive with cache & Iterative, fill table \\
\rowcolor{zebra}
Direction & Problem \(\rightarrow\) subproblems & Base cases \(\rightarrow\) problem \\
Space & O(depth) stack + memo & O(states) table \\
\rowcolor{zebra}
Complexity & Same asymptotically & Often faster constants \\
Ease & Easier to implement & Better for optimization \\
\rowcolor{zebra}
When to use & Not all subproblems needed & All subproblems computed \\
\end{longtable}

\setcodelang{Python}

\begin{Shaded}
\begin{Highlighting}[]
\CommentTok{\# Memoization (Fibonacci):}
\ImportTok{from}\NormalTok{ functools }\ImportTok{import}\NormalTok{ lru\_cache}

\AttributeTok{@lru\_cache}\NormalTok{(maxsize}\OperatorTok{=}\VariableTok{None}\NormalTok{)}
\KeywordTok{def}\NormalTok{ fib\_memo(n):}
    \ControlFlowTok{if}\NormalTok{ n }\OperatorTok{\textless{}=} \DecValTok{1}\NormalTok{: }\ControlFlowTok{return}\NormalTok{ n}
    \ControlFlowTok{return}\NormalTok{ fib\_memo(n}\OperatorTok{{-}}\DecValTok{1}\NormalTok{) }\OperatorTok{+}\NormalTok{ fib\_memo(n}\OperatorTok{{-}}\DecValTok{2}\NormalTok{)}

\CommentTok{\# Tabulation (Fibonacci):}
\KeywordTok{def}\NormalTok{ fib\_tab(n):}
    \ControlFlowTok{if}\NormalTok{ n }\OperatorTok{\textless{}=} \DecValTok{1}\NormalTok{: }\ControlFlowTok{return}\NormalTok{ n}
\NormalTok{    dp }\OperatorTok{=}\NormalTok{ [}\DecValTok{0}\NormalTok{, }\DecValTok{1}\NormalTok{]}
    \ControlFlowTok{for}\NormalTok{ i }\KeywordTok{in} \BuiltInTok{range}\NormalTok{(}\DecValTok{2}\NormalTok{, n}\OperatorTok{+}\DecValTok{1}\NormalTok{):}
\NormalTok{        dp.append(dp[}\OperatorTok{{-}}\DecValTok{1}\NormalTok{] }\OperatorTok{+}\NormalTok{ dp[}\OperatorTok{{-}}\DecValTok{2}\NormalTok{])}
    \ControlFlowTok{return}\NormalTok{ dp[n]}

\CommentTok{\# Space{-}optimized (only need last 2):}
\KeywordTok{def}\NormalTok{ fib\_opt(n):}
\NormalTok{    a, b }\OperatorTok{=} \DecValTok{0}\NormalTok{, }\DecValTok{1}
    \ControlFlowTok{for}\NormalTok{ \_ }\KeywordTok{in} \BuiltInTok{range}\NormalTok{(n):}
\NormalTok{        a, b }\OperatorTok{=}\NormalTok{ b, a }\OperatorTok{+}\NormalTok{ b}
    \ControlFlowTok{return}\NormalTok{ a}
\end{Highlighting}
\end{Shaded}

\subsubsection{How to Identify DP Problems}\label{how-to-identify-dp-problems}

\textbf{DP applies when:}

\begin{enumerate}
\def\labelenumi{\arabic{enumi}.}
\tightlist
\item
  Problem asks for \textbf{optimal value} (min/max) or \textbf{count of ways}
\item
  Subproblems \textbf{overlap} (same subproblem solved repeatedly)
\item
  Problem has \textbf{optimal substructure} (optimal solution uses optimal solutions to subproblems)
\item
  Problem asks: "is it possible to achieve X?"
\end{enumerate}

\textbf{DP does NOT apply when:}

\begin{itemize}
\tightlist
\item
  No overlapping subproblems (use divide and conquer)
\item
  Greedy works (exchange argument)
\item
  Subproblems are independent
\end{itemize}

\textbf{DP design recipe:}

\begin{enumerate}
\def\labelenumi{\arabic{enumi}.}
\tightlist
\item
  Define the \textbf{state} --- what information uniquely identifies a subproblem?
\item
  Write the \textbf{recurrence} (transition function)
\item
  Identify \textbf{base cases}
\item
  Determine \textbf{computation order} (which states depend on which)
\item
  Answer --- which state(s) give the final answer?
\end{enumerate}

\subsubsection{Classic DP Families}\label{classic-dp-families}

\textbf{0/1 Knapsack:}
Items with weight and value; capacity W; each item used at most once.
Recurrence: \texttt{dp{[}i{]}{[}w{]}\ =\ max(dp{[}i-1{]}{[}w{]},\ dp{[}i-1{]}{[}w-weight{[}i{]}{]}\ +\ value{[}i{]})}

\setcodelang{Python}

\begin{Shaded}
\begin{Highlighting}[]
\KeywordTok{def}\NormalTok{ knapsack\_01(weights, values, W):}
\NormalTok{    n }\OperatorTok{=} \BuiltInTok{len}\NormalTok{(weights)}
\NormalTok{    dp }\OperatorTok{=}\NormalTok{ [[}\DecValTok{0}\NormalTok{] }\OperatorTok{*}\NormalTok{ (W }\OperatorTok{+} \DecValTok{1}\NormalTok{) }\ControlFlowTok{for}\NormalTok{ \_ }\KeywordTok{in} \BuiltInTok{range}\NormalTok{(n }\OperatorTok{+} \DecValTok{1}\NormalTok{)]}
    \ControlFlowTok{for}\NormalTok{ i }\KeywordTok{in} \BuiltInTok{range}\NormalTok{(}\DecValTok{1}\NormalTok{, n }\OperatorTok{+} \DecValTok{1}\NormalTok{):}
        \ControlFlowTok{for}\NormalTok{ w }\KeywordTok{in} \BuiltInTok{range}\NormalTok{(W }\OperatorTok{+} \DecValTok{1}\NormalTok{):}
\NormalTok{            dp[i][w] }\OperatorTok{=}\NormalTok{ dp[i}\OperatorTok{{-}}\DecValTok{1}\NormalTok{][w]   }\CommentTok{\# don\textquotesingle{}t take item i}
            \ControlFlowTok{if}\NormalTok{ weights[i}\OperatorTok{{-}}\DecValTok{1}\NormalTok{] }\OperatorTok{\textless{}=}\NormalTok{ w:}
\NormalTok{                dp[i][w] }\OperatorTok{=} \BuiltInTok{max}\NormalTok{(dp[i][w], dp[i}\OperatorTok{{-}}\DecValTok{1}\NormalTok{][w }\OperatorTok{{-}}\NormalTok{ weights[i}\OperatorTok{{-}}\DecValTok{1}\NormalTok{]] }\OperatorTok{+}\NormalTok{ values[i}\OperatorTok{{-}}\DecValTok{1}\NormalTok{])}
    \ControlFlowTok{return}\NormalTok{ dp[n][W]}

\CommentTok{\# Space{-}optimized (1D DP) — iterate w in reverse:}
\KeywordTok{def}\NormalTok{ knapsack\_01\_opt(weights, values, W):}
\NormalTok{    dp }\OperatorTok{=}\NormalTok{ [}\DecValTok{0}\NormalTok{] }\OperatorTok{*}\NormalTok{ (W }\OperatorTok{+} \DecValTok{1}\NormalTok{)}
    \ControlFlowTok{for}\NormalTok{ i }\KeywordTok{in} \BuiltInTok{range}\NormalTok{(}\BuiltInTok{len}\NormalTok{(weights)):}
        \ControlFlowTok{for}\NormalTok{ w }\KeywordTok{in} \BuiltInTok{range}\NormalTok{(W, weights[i] }\OperatorTok{{-}} \DecValTok{1}\NormalTok{, }\OperatorTok{{-}}\DecValTok{1}\NormalTok{):   }\CommentTok{\# MUST go right to left}
\NormalTok{            dp[w] }\OperatorTok{=} \BuiltInTok{max}\NormalTok{(dp[w], dp[w }\OperatorTok{{-}}\NormalTok{ weights[i]] }\OperatorTok{+}\NormalTok{ values[i])}
    \ControlFlowTok{return}\NormalTok{ dp[W]}
\end{Highlighting}
\end{Shaded}

\textbf{Unbounded Knapsack (items can be reused):}

\setcodelang{Python}

\begin{Shaded}
\begin{Highlighting}[]
\KeywordTok{def}\NormalTok{ knapsack\_unbounded(weights, values, W):}
\NormalTok{    dp }\OperatorTok{=}\NormalTok{ [}\DecValTok{0}\NormalTok{] }\OperatorTok{*}\NormalTok{ (W }\OperatorTok{+} \DecValTok{1}\NormalTok{)}
    \ControlFlowTok{for}\NormalTok{ w }\KeywordTok{in} \BuiltInTok{range}\NormalTok{(W }\OperatorTok{+} \DecValTok{1}\NormalTok{):}
        \ControlFlowTok{for}\NormalTok{ i }\KeywordTok{in} \BuiltInTok{range}\NormalTok{(}\BuiltInTok{len}\NormalTok{(weights)):}
            \ControlFlowTok{if}\NormalTok{ weights[i] }\OperatorTok{\textless{}=}\NormalTok{ w:}
\NormalTok{                dp[w] }\OperatorTok{=} \BuiltInTok{max}\NormalTok{(dp[w], dp[w }\OperatorTok{{-}}\NormalTok{ weights[i]] }\OperatorTok{+}\NormalTok{ values[i])}
    \ControlFlowTok{return}\NormalTok{ dp[W]}
\end{Highlighting}
\end{Shaded}

\textbf{Coin Change (minimum coins):}
Recurrence: \texttt{dp{[}amount{]}\ =\ min(dp{[}amount\ -\ coin{]}\ +\ 1)\ for\ each\ coin}

\setcodelang{Python}

\begin{Shaded}
\begin{Highlighting}[]
\KeywordTok{def}\NormalTok{ coin\_change(coins, amount):}
\NormalTok{    dp }\OperatorTok{=}\NormalTok{ [}\BuiltInTok{float}\NormalTok{(}\StringTok{\textquotesingle{}inf\textquotesingle{}}\NormalTok{)] }\OperatorTok{*}\NormalTok{ (amount }\OperatorTok{+} \DecValTok{1}\NormalTok{)}
\NormalTok{    dp[}\DecValTok{0}\NormalTok{] }\OperatorTok{=} \DecValTok{0}
    \ControlFlowTok{for}\NormalTok{ a }\KeywordTok{in} \BuiltInTok{range}\NormalTok{(}\DecValTok{1}\NormalTok{, amount }\OperatorTok{+} \DecValTok{1}\NormalTok{):}
        \ControlFlowTok{for}\NormalTok{ coin }\KeywordTok{in}\NormalTok{ coins:}
            \ControlFlowTok{if}\NormalTok{ coin }\OperatorTok{\textless{}=}\NormalTok{ a:}
\NormalTok{                dp[a] }\OperatorTok{=} \BuiltInTok{min}\NormalTok{(dp[a], dp[a }\OperatorTok{{-}}\NormalTok{ coin] }\OperatorTok{+} \DecValTok{1}\NormalTok{)}
    \ControlFlowTok{return}\NormalTok{ dp[amount] }\ControlFlowTok{if}\NormalTok{ dp[amount] }\OperatorTok{!=} \BuiltInTok{float}\NormalTok{(}\StringTok{\textquotesingle{}inf\textquotesingle{}}\NormalTok{) }\ControlFlowTok{else} \OperatorTok{{-}}\DecValTok{1}
\end{Highlighting}
\end{Shaded}

\textbf{Coin Change (number of ways):}

\setcodelang{Python}

\begin{Shaded}
\begin{Highlighting}[]
\KeywordTok{def}\NormalTok{ coin\_change\_ways(coins, amount):}
\NormalTok{    dp }\OperatorTok{=}\NormalTok{ [}\DecValTok{0}\NormalTok{] }\OperatorTok{*}\NormalTok{ (amount }\OperatorTok{+} \DecValTok{1}\NormalTok{)}
\NormalTok{    dp[}\DecValTok{0}\NormalTok{] }\OperatorTok{=} \DecValTok{1}
    \ControlFlowTok{for}\NormalTok{ coin }\KeywordTok{in}\NormalTok{ coins:          }\CommentTok{\# order matters: outer=coins for combinations}
        \ControlFlowTok{for}\NormalTok{ a }\KeywordTok{in} \BuiltInTok{range}\NormalTok{(coin, amount }\OperatorTok{+} \DecValTok{1}\NormalTok{):}
\NormalTok{            dp[a] }\OperatorTok{+=}\NormalTok{ dp[a }\OperatorTok{{-}}\NormalTok{ coin]}
    \ControlFlowTok{return}\NormalTok{ dp[amount]}
\end{Highlighting}
\end{Shaded}

\textbf{Longest Common Subsequence (LCS):}
Recurrence:

\begin{itemize}
\tightlist
\item
  If s1{[}i{]} == s2{[}j{]}: \texttt{dp{[}i{]}{[}j{]}\ =\ dp{[}i-1{]}{[}j-1{]}\ +\ 1}
\item
  Else: \texttt{dp{[}i{]}{[}j{]}\ =\ max(dp{[}i-1{]}{[}j{]},\ dp{[}i{]}{[}j-1{]})}
\end{itemize}

\setcodelang{Python}

\begin{Shaded}
\begin{Highlighting}[]
\KeywordTok{def}\NormalTok{ lcs(s1, s2):}
\NormalTok{    m, n }\OperatorTok{=} \BuiltInTok{len}\NormalTok{(s1), }\BuiltInTok{len}\NormalTok{(s2)}
\NormalTok{    dp }\OperatorTok{=}\NormalTok{ [[}\DecValTok{0}\NormalTok{] }\OperatorTok{*}\NormalTok{ (n }\OperatorTok{+} \DecValTok{1}\NormalTok{) }\ControlFlowTok{for}\NormalTok{ \_ }\KeywordTok{in} \BuiltInTok{range}\NormalTok{(m }\OperatorTok{+} \DecValTok{1}\NormalTok{)]}
    \ControlFlowTok{for}\NormalTok{ i }\KeywordTok{in} \BuiltInTok{range}\NormalTok{(}\DecValTok{1}\NormalTok{, m }\OperatorTok{+} \DecValTok{1}\NormalTok{):}
        \ControlFlowTok{for}\NormalTok{ j }\KeywordTok{in} \BuiltInTok{range}\NormalTok{(}\DecValTok{1}\NormalTok{, n }\OperatorTok{+} \DecValTok{1}\NormalTok{):}
            \ControlFlowTok{if}\NormalTok{ s1[i}\OperatorTok{{-}}\DecValTok{1}\NormalTok{] }\OperatorTok{==}\NormalTok{ s2[j}\OperatorTok{{-}}\DecValTok{1}\NormalTok{]:}
\NormalTok{                dp[i][j] }\OperatorTok{=}\NormalTok{ dp[i}\OperatorTok{{-}}\DecValTok{1}\NormalTok{][j}\OperatorTok{{-}}\DecValTok{1}\NormalTok{] }\OperatorTok{+} \DecValTok{1}
            \ControlFlowTok{else}\NormalTok{:}
\NormalTok{                dp[i][j] }\OperatorTok{=} \BuiltInTok{max}\NormalTok{(dp[i}\OperatorTok{{-}}\DecValTok{1}\NormalTok{][j], dp[i][j}\OperatorTok{{-}}\DecValTok{1}\NormalTok{])}
    \ControlFlowTok{return}\NormalTok{ dp[m][n]}
\end{Highlighting}
\end{Shaded}

\textbf{Edit Distance (Levenshtein):}
Operations: insert, delete, replace (each cost 1).
Recurrence:

\begin{itemize}
\tightlist
\item
  If s1{[}i{]} == s2{[}j{]}: \texttt{dp{[}i{]}{[}j{]}\ =\ dp{[}i-1{]}{[}j-1{]}}
\item
  Else: \texttt{dp{[}i{]}{[}j{]}\ =\ 1\ +\ min(dp{[}i-1{]}{[}j{]},\ dp{[}i{]}{[}j-1{]},\ dp{[}i-1{]}{[}j-1{]})}
\end{itemize}

\setcodelang{Python}

\begin{Shaded}
\begin{Highlighting}[]
\KeywordTok{def}\NormalTok{ edit\_distance(s1, s2):}
\NormalTok{    m, n }\OperatorTok{=} \BuiltInTok{len}\NormalTok{(s1), }\BuiltInTok{len}\NormalTok{(s2)}
\NormalTok{    dp }\OperatorTok{=}\NormalTok{ [[}\DecValTok{0}\NormalTok{] }\OperatorTok{*}\NormalTok{ (n }\OperatorTok{+} \DecValTok{1}\NormalTok{) }\ControlFlowTok{for}\NormalTok{ \_ }\KeywordTok{in} \BuiltInTok{range}\NormalTok{(m }\OperatorTok{+} \DecValTok{1}\NormalTok{)]}
    \ControlFlowTok{for}\NormalTok{ i }\KeywordTok{in} \BuiltInTok{range}\NormalTok{(m }\OperatorTok{+} \DecValTok{1}\NormalTok{): dp[i][}\DecValTok{0}\NormalTok{] }\OperatorTok{=}\NormalTok{ i}
    \ControlFlowTok{for}\NormalTok{ j }\KeywordTok{in} \BuiltInTok{range}\NormalTok{(n }\OperatorTok{+} \DecValTok{1}\NormalTok{): dp[}\DecValTok{0}\NormalTok{][j] }\OperatorTok{=}\NormalTok{ j}
    \ControlFlowTok{for}\NormalTok{ i }\KeywordTok{in} \BuiltInTok{range}\NormalTok{(}\DecValTok{1}\NormalTok{, m }\OperatorTok{+} \DecValTok{1}\NormalTok{):}
        \ControlFlowTok{for}\NormalTok{ j }\KeywordTok{in} \BuiltInTok{range}\NormalTok{(}\DecValTok{1}\NormalTok{, n }\OperatorTok{+} \DecValTok{1}\NormalTok{):}
            \ControlFlowTok{if}\NormalTok{ s1[i}\OperatorTok{{-}}\DecValTok{1}\NormalTok{] }\OperatorTok{==}\NormalTok{ s2[j}\OperatorTok{{-}}\DecValTok{1}\NormalTok{]:}
\NormalTok{                dp[i][j] }\OperatorTok{=}\NormalTok{ dp[i}\OperatorTok{{-}}\DecValTok{1}\NormalTok{][j}\OperatorTok{{-}}\DecValTok{1}\NormalTok{]}
            \ControlFlowTok{else}\NormalTok{:}
\NormalTok{                dp[i][j] }\OperatorTok{=} \DecValTok{1} \OperatorTok{+} \BuiltInTok{min}\NormalTok{(dp[i}\OperatorTok{{-}}\DecValTok{1}\NormalTok{][j],     }\CommentTok{\# delete from s1}
\NormalTok{                                   dp[i][j}\OperatorTok{{-}}\DecValTok{1}\NormalTok{],      }\CommentTok{\# insert into s1}
\NormalTok{                                   dp[i}\OperatorTok{{-}}\DecValTok{1}\NormalTok{][j}\OperatorTok{{-}}\DecValTok{1}\NormalTok{])    }\CommentTok{\# replace}
    \ControlFlowTok{return}\NormalTok{ dp[m][n]}
\end{Highlighting}
\end{Shaded}

\textbf{Longest Increasing Subsequence (LIS):}

\setcodelang{Python}

\begin{Shaded}
\begin{Highlighting}[]
\CommentTok{\# O(n²) DP:}
\KeywordTok{def}\NormalTok{ lis\_n2(arr):}
\NormalTok{    n }\OperatorTok{=} \BuiltInTok{len}\NormalTok{(arr)}
\NormalTok{    dp }\OperatorTok{=}\NormalTok{ [}\DecValTok{1}\NormalTok{] }\OperatorTok{*}\NormalTok{ n}
    \ControlFlowTok{for}\NormalTok{ i }\KeywordTok{in} \BuiltInTok{range}\NormalTok{(}\DecValTok{1}\NormalTok{, n):}
        \ControlFlowTok{for}\NormalTok{ j }\KeywordTok{in} \BuiltInTok{range}\NormalTok{(i):}
            \ControlFlowTok{if}\NormalTok{ arr[j] }\OperatorTok{\textless{}}\NormalTok{ arr[i]:}
\NormalTok{                dp[i] }\OperatorTok{=} \BuiltInTok{max}\NormalTok{(dp[i], dp[j] }\OperatorTok{+} \DecValTok{1}\NormalTok{)}
    \ControlFlowTok{return} \BuiltInTok{max}\NormalTok{(dp)}

\CommentTok{\# O(n log n) patience sorting:}
\ImportTok{import}\NormalTok{ bisect}

\KeywordTok{def}\NormalTok{ lis\_nlogn(arr):}
\NormalTok{    tails }\OperatorTok{=}\NormalTok{ []  }\CommentTok{\# tails[i] = smallest tail element of all IS of length i+1}
    \ControlFlowTok{for}\NormalTok{ x }\KeywordTok{in}\NormalTok{ arr:}
\NormalTok{        pos }\OperatorTok{=}\NormalTok{ bisect.bisect\_left(tails, x)}
        \ControlFlowTok{if}\NormalTok{ pos }\OperatorTok{==} \BuiltInTok{len}\NormalTok{(tails):}
\NormalTok{            tails.append(x)}
        \ControlFlowTok{else}\NormalTok{:}
\NormalTok{            tails[pos] }\OperatorTok{=}\NormalTok{ x}
    \ControlFlowTok{return} \BuiltInTok{len}\NormalTok{(tails)}
\end{Highlighting}
\end{Shaded}

\textbf{Matrix Chain Multiplication (Interval DP):}
Recurrence: \texttt{dp{[}i{]}{[}j{]}\ =\ min\ over\ k\ of\ (dp{[}i{]}{[}k{]}\ +\ dp{[}k+1{]}{[}j{]}\ +\ dims{[}i-1{]}*dims{[}k{]}*dims{[}j{]})}

\setcodelang{Python}

\begin{Shaded}
\begin{Highlighting}[]
\KeywordTok{def}\NormalTok{ matrix\_chain(dims):}
\NormalTok{    n }\OperatorTok{=} \BuiltInTok{len}\NormalTok{(dims) }\OperatorTok{{-}} \DecValTok{1}   \CommentTok{\# number of matrices}
\NormalTok{    dp }\OperatorTok{=}\NormalTok{ [[}\DecValTok{0}\NormalTok{] }\OperatorTok{*}\NormalTok{ n }\ControlFlowTok{for}\NormalTok{ \_ }\KeywordTok{in} \BuiltInTok{range}\NormalTok{(n)]}
    \ControlFlowTok{for}\NormalTok{ length }\KeywordTok{in} \BuiltInTok{range}\NormalTok{(}\DecValTok{2}\NormalTok{, n }\OperatorTok{+} \DecValTok{1}\NormalTok{):      }\CommentTok{\# subproblem length}
        \ControlFlowTok{for}\NormalTok{ i }\KeywordTok{in} \BuiltInTok{range}\NormalTok{(n }\OperatorTok{{-}}\NormalTok{ length }\OperatorTok{+} \DecValTok{1}\NormalTok{):}
\NormalTok{            j }\OperatorTok{=}\NormalTok{ i }\OperatorTok{+}\NormalTok{ length }\OperatorTok{{-}} \DecValTok{1}
\NormalTok{            dp[i][j] }\OperatorTok{=} \BuiltInTok{float}\NormalTok{(}\StringTok{\textquotesingle{}inf\textquotesingle{}}\NormalTok{)}
            \ControlFlowTok{for}\NormalTok{ k }\KeywordTok{in} \BuiltInTok{range}\NormalTok{(i, j):}
\NormalTok{                cost }\OperatorTok{=}\NormalTok{ dp[i][k] }\OperatorTok{+}\NormalTok{ dp[k}\OperatorTok{+}\DecValTok{1}\NormalTok{][j] }\OperatorTok{+}\NormalTok{ dims[i] }\OperatorTok{*}\NormalTok{ dims[k}\OperatorTok{+}\DecValTok{1}\NormalTok{] }\OperatorTok{*}\NormalTok{ dims[j}\OperatorTok{+}\DecValTok{1}\NormalTok{]}
\NormalTok{                dp[i][j] }\OperatorTok{=} \BuiltInTok{min}\NormalTok{(dp[i][j], cost)}
    \ControlFlowTok{return}\NormalTok{ dp[}\DecValTok{0}\NormalTok{][n}\OperatorTok{{-}}\DecValTok{1}\NormalTok{]}
\end{Highlighting}
\end{Shaded}

\textbf{Matrix Paths (unique paths, min cost path):}

\setcodelang{Python}

\begin{Shaded}
\begin{Highlighting}[]
\KeywordTok{def}\NormalTok{ unique\_paths(m, n):}
\NormalTok{    dp }\OperatorTok{=}\NormalTok{ [[}\DecValTok{1}\NormalTok{] }\OperatorTok{*}\NormalTok{ n }\ControlFlowTok{for}\NormalTok{ \_ }\KeywordTok{in} \BuiltInTok{range}\NormalTok{(m)]}
    \ControlFlowTok{for}\NormalTok{ i }\KeywordTok{in} \BuiltInTok{range}\NormalTok{(}\DecValTok{1}\NormalTok{, m):}
        \ControlFlowTok{for}\NormalTok{ j }\KeywordTok{in} \BuiltInTok{range}\NormalTok{(}\DecValTok{1}\NormalTok{, n):}
\NormalTok{            dp[i][j] }\OperatorTok{=}\NormalTok{ dp[i}\OperatorTok{{-}}\DecValTok{1}\NormalTok{][j] }\OperatorTok{+}\NormalTok{ dp[i][j}\OperatorTok{{-}}\DecValTok{1}\NormalTok{]}
    \ControlFlowTok{return}\NormalTok{ dp[m}\OperatorTok{{-}}\DecValTok{1}\NormalTok{][n}\OperatorTok{{-}}\DecValTok{1}\NormalTok{]}

\KeywordTok{def}\NormalTok{ min\_path\_sum(grid):}
\NormalTok{    m, n }\OperatorTok{=} \BuiltInTok{len}\NormalTok{(grid), }\BuiltInTok{len}\NormalTok{(grid[}\DecValTok{0}\NormalTok{])}
\NormalTok{    dp }\OperatorTok{=}\NormalTok{ [[}\DecValTok{0}\NormalTok{] }\OperatorTok{*}\NormalTok{ n }\ControlFlowTok{for}\NormalTok{ \_ }\KeywordTok{in} \BuiltInTok{range}\NormalTok{(m)]}
\NormalTok{    dp[}\DecValTok{0}\NormalTok{][}\DecValTok{0}\NormalTok{] }\OperatorTok{=}\NormalTok{ grid[}\DecValTok{0}\NormalTok{][}\DecValTok{0}\NormalTok{]}
    \ControlFlowTok{for}\NormalTok{ i }\KeywordTok{in} \BuiltInTok{range}\NormalTok{(}\DecValTok{1}\NormalTok{, m): dp[i][}\DecValTok{0}\NormalTok{] }\OperatorTok{=}\NormalTok{ dp[i}\OperatorTok{{-}}\DecValTok{1}\NormalTok{][}\DecValTok{0}\NormalTok{] }\OperatorTok{+}\NormalTok{ grid[i][}\DecValTok{0}\NormalTok{]}
    \ControlFlowTok{for}\NormalTok{ j }\KeywordTok{in} \BuiltInTok{range}\NormalTok{(}\DecValTok{1}\NormalTok{, n): dp[}\DecValTok{0}\NormalTok{][j] }\OperatorTok{=}\NormalTok{ dp[}\DecValTok{0}\NormalTok{][j}\OperatorTok{{-}}\DecValTok{1}\NormalTok{] }\OperatorTok{+}\NormalTok{ grid[}\DecValTok{0}\NormalTok{][j]}
    \ControlFlowTok{for}\NormalTok{ i }\KeywordTok{in} \BuiltInTok{range}\NormalTok{(}\DecValTok{1}\NormalTok{, m):}
        \ControlFlowTok{for}\NormalTok{ j }\KeywordTok{in} \BuiltInTok{range}\NormalTok{(}\DecValTok{1}\NormalTok{, n):}
\NormalTok{            dp[i][j] }\OperatorTok{=}\NormalTok{ grid[i][j] }\OperatorTok{+} \BuiltInTok{min}\NormalTok{(dp[i}\OperatorTok{{-}}\DecValTok{1}\NormalTok{][j], dp[i][j}\OperatorTok{{-}}\DecValTok{1}\NormalTok{])}
    \ControlFlowTok{return}\NormalTok{ dp[m}\OperatorTok{{-}}\DecValTok{1}\NormalTok{][n}\OperatorTok{{-}}\DecValTok{1}\NormalTok{]}
\end{Highlighting}
\end{Shaded}

\textbf{DP on Trees:}

\setcodelang{Python}

\begin{Shaded}
\begin{Highlighting}[]
\KeywordTok{def}\NormalTok{ tree\_dp\_max\_independent\_set(adj, n):}
    \CommentTok{\# dp[node][0] = max IS if node NOT included}
    \CommentTok{\# dp[node][1] = max IS if node included}
\NormalTok{    dp }\OperatorTok{=}\NormalTok{ [[}\DecValTok{0}\NormalTok{, }\DecValTok{1}\NormalTok{] }\ControlFlowTok{for}\NormalTok{ \_ }\KeywordTok{in} \BuiltInTok{range}\NormalTok{(n)]   }\CommentTok{\# [0] = not taken, [1] = taken (value 1)}

    \KeywordTok{def}\NormalTok{ dfs(u, parent):}
        \ControlFlowTok{for}\NormalTok{ v }\KeywordTok{in}\NormalTok{ adj[u]:}
            \ControlFlowTok{if}\NormalTok{ v }\OperatorTok{==}\NormalTok{ parent: }\ControlFlowTok{continue}
\NormalTok{            dfs(v, u)}
            \CommentTok{\# If u not taken, v can be taken or not}
\NormalTok{            dp[u][}\DecValTok{0}\NormalTok{] }\OperatorTok{+=} \BuiltInTok{max}\NormalTok{(dp[v][}\DecValTok{0}\NormalTok{], dp[v][}\DecValTok{1}\NormalTok{])}
            \CommentTok{\# If u taken, v cannot be taken}
\NormalTok{            dp[u][}\DecValTok{1}\NormalTok{] }\OperatorTok{+=}\NormalTok{ dp[v][}\DecValTok{0}\NormalTok{]}

\NormalTok{    dfs(}\DecValTok{0}\NormalTok{, }\OperatorTok{{-}}\DecValTok{1}\NormalTok{)}
    \ControlFlowTok{return} \BuiltInTok{max}\NormalTok{(dp[}\DecValTok{0}\NormalTok{][}\DecValTok{0}\NormalTok{], dp[}\DecValTok{0}\NormalTok{][}\DecValTok{1}\NormalTok{])}
\end{Highlighting}
\end{Shaded}

\textbf{Bitmask DP (Travelling Salesman):}
State: \texttt{dp{[}mask{]}{[}i{]}} = minimum cost to visit all cities in \texttt{mask}, ending at city \texttt{i}.

\setcodelang{Python}

\begin{Shaded}
\begin{Highlighting}[]
\KeywordTok{def}\NormalTok{ tsp(dist, n):}
\NormalTok{    INF }\OperatorTok{=} \BuiltInTok{float}\NormalTok{(}\StringTok{\textquotesingle{}inf\textquotesingle{}}\NormalTok{)}
\NormalTok{    dp }\OperatorTok{=}\NormalTok{ [[INF] }\OperatorTok{*}\NormalTok{ n }\ControlFlowTok{for}\NormalTok{ \_ }\KeywordTok{in} \BuiltInTok{range}\NormalTok{(}\DecValTok{1} \OperatorTok{\textless{}\textless{}}\NormalTok{ n)]}
\NormalTok{    dp[}\DecValTok{1}\NormalTok{][}\DecValTok{0}\NormalTok{] }\OperatorTok{=} \DecValTok{0}   \CommentTok{\# visited city 0, at city 0}

    \ControlFlowTok{for}\NormalTok{ mask }\KeywordTok{in} \BuiltInTok{range}\NormalTok{(}\DecValTok{1}\NormalTok{, }\DecValTok{1} \OperatorTok{\textless{}\textless{}}\NormalTok{ n):}
        \ControlFlowTok{for}\NormalTok{ u }\KeywordTok{in} \BuiltInTok{range}\NormalTok{(n):}
            \ControlFlowTok{if} \KeywordTok{not}\NormalTok{ (mask }\OperatorTok{\textgreater{}\textgreater{}}\NormalTok{ u }\OperatorTok{\&} \DecValTok{1}\NormalTok{): }\ControlFlowTok{continue}
            \ControlFlowTok{if}\NormalTok{ dp[mask][u] }\OperatorTok{==}\NormalTok{ INF: }\ControlFlowTok{continue}
            \ControlFlowTok{for}\NormalTok{ v }\KeywordTok{in} \BuiltInTok{range}\NormalTok{(n):}
                \ControlFlowTok{if}\NormalTok{ mask }\OperatorTok{\textgreater{}\textgreater{}}\NormalTok{ v }\OperatorTok{\&} \DecValTok{1}\NormalTok{: }\ControlFlowTok{continue}   \CommentTok{\# already visited}
\NormalTok{                new\_mask }\OperatorTok{=}\NormalTok{ mask }\OperatorTok{|}\NormalTok{ (}\DecValTok{1} \OperatorTok{\textless{}\textless{}}\NormalTok{ v)}
\NormalTok{                dp[new\_mask][v] }\OperatorTok{=} \BuiltInTok{min}\NormalTok{(dp[new\_mask][v], dp[mask][u] }\OperatorTok{+}\NormalTok{ dist[u][v])}

\NormalTok{    full }\OperatorTok{=}\NormalTok{ (}\DecValTok{1} \OperatorTok{\textless{}\textless{}}\NormalTok{ n) }\OperatorTok{{-}} \DecValTok{1}
    \ControlFlowTok{return} \BuiltInTok{min}\NormalTok{(dp[full][v] }\OperatorTok{+}\NormalTok{ dist[v][}\DecValTok{0}\NormalTok{] }\ControlFlowTok{for}\NormalTok{ v }\KeywordTok{in} \BuiltInTok{range}\NormalTok{(n))}
\end{Highlighting}
\end{Shaded}

\textbf{DP Key Takeaways:}

\begin{itemize}
\tightlist
\item
  \textbf{State is everything} --- spend most time defining what uniquely describes a subproblem
\item
  \textbf{Recurrence is mechanical} --- once state is defined, transitions follow naturally
\item
  \textbf{Base cases are critical} --- wrong base cases cause wrong answers even with correct recurrence
\item
  \textbf{Space optimization:} 2D DP can often become 1D; 1D can become O(1) if only adjacent states needed
\end{itemize}

\subsection{Greedy Algorithms}\label{greedy-algorithms}

Greedy algorithms make the \textbf{locally optimal choice at each step} hoping to reach a global optimum.

\subsubsection{When Greedy Works}\label{when-greedy-works}

Greedy is correct when two conditions hold:

\begin{enumerate}
\def\labelenumi{\arabic{enumi}.}
\tightlist
\item
  \textbf{Greedy choice property:} A globally optimal solution can always be reached by making locally optimal choices
\item
  \textbf{Optimal substructure:} Optimal solution contains optimal solutions to subproblems
\end{enumerate}

\textbf{Exchange argument:} Standard proof technique. Assume greedy is not optimal. Show that any optimal solution can be transformed (by swapping) to match greedy without increasing cost --- contradiction.

\subsubsection{Classic Greedy Problems}\label{classic-greedy-problems}

\textbf{Activity Selection (Interval Scheduling):}
Select maximum number of non-overlapping intervals.
Strategy: Sort by end time; always pick the activity that finishes earliest.

\setcodelang{Python}

\begin{Shaded}
\begin{Highlighting}[]
\KeywordTok{def}\NormalTok{ activity\_selection(intervals):}
\NormalTok{    intervals.sort(key}\OperatorTok{=}\KeywordTok{lambda}\NormalTok{ x: x[}\DecValTok{1}\NormalTok{])   }\CommentTok{\# sort by end time}
\NormalTok{    count }\OperatorTok{=} \DecValTok{0}
\NormalTok{    last\_end }\OperatorTok{=} \BuiltInTok{float}\NormalTok{(}\StringTok{\textquotesingle{}{-}inf\textquotesingle{}}\NormalTok{)}
    \ControlFlowTok{for}\NormalTok{ start, end }\KeywordTok{in}\NormalTok{ intervals:}
        \ControlFlowTok{if}\NormalTok{ start }\OperatorTok{\textgreater{}=}\NormalTok{ last\_end:}
\NormalTok{            count }\OperatorTok{+=} \DecValTok{1}
\NormalTok{            last\_end }\OperatorTok{=}\NormalTok{ end}
    \ControlFlowTok{return}\NormalTok{ count}
\end{Highlighting}
\end{Shaded}

\textbf{Huffman Coding:} Build optimal prefix-free codes using a min-heap. Greedy: always merge the two lowest-frequency nodes.

\textbf{Fractional Knapsack:} (unlike 0/1 knapsack, greedy works here)

\setcodelang{Python}

\begin{Shaded}
\begin{Highlighting}[]
\KeywordTok{def}\NormalTok{ fractional\_knapsack(weights, values, W):}
\NormalTok{    ratios }\OperatorTok{=} \BuiltInTok{sorted}\NormalTok{(}\BuiltInTok{zip}\NormalTok{(values, weights), key}\OperatorTok{=}\KeywordTok{lambda}\NormalTok{ x: x[}\DecValTok{0}\NormalTok{]}\OperatorTok{/}\NormalTok{x[}\DecValTok{1}\NormalTok{], reverse}\OperatorTok{=}\VariableTok{True}\NormalTok{)}
\NormalTok{    total }\OperatorTok{=} \DecValTok{0}
    \ControlFlowTok{for}\NormalTok{ v, w }\KeywordTok{in}\NormalTok{ ratios:}
        \ControlFlowTok{if}\NormalTok{ W }\OperatorTok{\textgreater{}=}\NormalTok{ w:}
\NormalTok{            total }\OperatorTok{+=}\NormalTok{ v}\OperatorTok{;}\NormalTok{ W }\OperatorTok{{-}=}\NormalTok{ w}
        \ControlFlowTok{else}\NormalTok{:}
\NormalTok{            total }\OperatorTok{+=}\NormalTok{ v }\OperatorTok{*}\NormalTok{ W }\OperatorTok{/}\NormalTok{ w}\OperatorTok{;} \ControlFlowTok{break}
    \ControlFlowTok{return}\NormalTok{ total}
\end{Highlighting}
\end{Shaded}

\textbf{Jump Game:} Can you reach the end of an array where each element is max jump length?

\setcodelang{Python}

\begin{Shaded}
\begin{Highlighting}[]
\KeywordTok{def}\NormalTok{ can\_jump(nums):}
\NormalTok{    max\_reach }\OperatorTok{=} \DecValTok{0}
    \ControlFlowTok{for}\NormalTok{ i, jump }\KeywordTok{in} \BuiltInTok{enumerate}\NormalTok{(nums):}
        \ControlFlowTok{if}\NormalTok{ i }\OperatorTok{\textgreater{}}\NormalTok{ max\_reach: }\ControlFlowTok{return} \VariableTok{False}
\NormalTok{        max\_reach }\OperatorTok{=} \BuiltInTok{max}\NormalTok{(max\_reach, i }\OperatorTok{+}\NormalTok{ jump)}
    \ControlFlowTok{return} \VariableTok{True}
\end{Highlighting}
\end{Shaded}

\textbf{Task Scheduling:} Schedule tasks with deadlines to maximize profit.
Strategy: Sort by profit descending; greedily assign tasks to latest available slot.

\textbf{Gas Station:} Find starting station to complete circular tour.

\setcodelang{Python}

\begin{Shaded}
\begin{Highlighting}[]
\KeywordTok{def}\NormalTok{ can\_complete\_circuit(gas, cost):}
\NormalTok{    total\_surplus }\OperatorTok{=} \DecValTok{0}
\NormalTok{    surplus }\OperatorTok{=} \DecValTok{0}
\NormalTok{    start }\OperatorTok{=} \DecValTok{0}
    \ControlFlowTok{for}\NormalTok{ i }\KeywordTok{in} \BuiltInTok{range}\NormalTok{(}\BuiltInTok{len}\NormalTok{(gas)):}
\NormalTok{        diff }\OperatorTok{=}\NormalTok{ gas[i] }\OperatorTok{{-}}\NormalTok{ cost[i]}
\NormalTok{        total\_surplus }\OperatorTok{+=}\NormalTok{ diff}
\NormalTok{        surplus }\OperatorTok{+=}\NormalTok{ diff}
        \ControlFlowTok{if}\NormalTok{ surplus }\OperatorTok{\textless{}} \DecValTok{0}\NormalTok{:}
\NormalTok{            start }\OperatorTok{=}\NormalTok{ i }\OperatorTok{+} \DecValTok{1}
\NormalTok{            surplus }\OperatorTok{=} \DecValTok{0}
    \ControlFlowTok{return}\NormalTok{ start }\ControlFlowTok{if}\NormalTok{ total\_surplus }\OperatorTok{\textgreater{}=} \DecValTok{0} \ControlFlowTok{else} \OperatorTok{{-}}\DecValTok{1}
\end{Highlighting}
\end{Shaded}

\textbf{Greedy vs DP:} Greedy is O(n log n) or O(n); DP is O(n\(^2\)) or worse. \textbf{Prefer greedy when provably correct} (exchange argument or known greedy structure). DP is the safe fallback.

\subsection{Recursion \& Backtracking}\label{recursion--backtracking}

\subsubsection{Recursion Template}\label{recursion-template}

\setcodelang{Python}

\begin{Shaded}
\begin{Highlighting}[]
\KeywordTok{def}\NormalTok{ solve(state, ...):}
    \CommentTok{\# Base case}
    \ControlFlowTok{if}\NormalTok{ is\_complete(state):}
\NormalTok{        record\_solution(state)}
        \ControlFlowTok{return}

    \CommentTok{\# Recursive case}
    \ControlFlowTok{for}\NormalTok{ choice }\KeywordTok{in}\NormalTok{ get\_choices(state):}
        \ControlFlowTok{if}\NormalTok{ is\_valid(choice, state):}
\NormalTok{            make\_choice(state, choice)}
\NormalTok{            solve(state, ...)}
\NormalTok{            undo\_choice(state, choice)   }\CommentTok{\# backtrack}
\end{Highlighting}
\end{Shaded}

\subsubsection{Backtracking Framework}\label{backtracking-framework}

Backtracking = DFS on the decision tree with \textbf{pruning} to avoid exploring invalid paths.

\setcodelang{Python}

\begin{Shaded}
\begin{Highlighting}[]
\CommentTok{\# Subsets (power set):}
\KeywordTok{def}\NormalTok{ subsets(nums):}
\NormalTok{    result }\OperatorTok{=}\NormalTok{ []}
    \KeywordTok{def}\NormalTok{ backtrack(start, path):}
\NormalTok{        result.append(path[:])}
        \ControlFlowTok{for}\NormalTok{ i }\KeywordTok{in} \BuiltInTok{range}\NormalTok{(start, }\BuiltInTok{len}\NormalTok{(nums)):}
\NormalTok{            path.append(nums[i])}
\NormalTok{            backtrack(i }\OperatorTok{+} \DecValTok{1}\NormalTok{, path)}
\NormalTok{            path.pop()}
\NormalTok{    backtrack(}\DecValTok{0}\NormalTok{, [])}
    \ControlFlowTok{return}\NormalTok{ result}

\CommentTok{\# Permutations:}
\KeywordTok{def}\NormalTok{ permutations(nums):}
\NormalTok{    result }\OperatorTok{=}\NormalTok{ []}
    \KeywordTok{def}\NormalTok{ backtrack(path, used):}
        \ControlFlowTok{if} \BuiltInTok{len}\NormalTok{(path) }\OperatorTok{==} \BuiltInTok{len}\NormalTok{(nums):}
\NormalTok{            result.append(path[:])}
            \ControlFlowTok{return}
        \ControlFlowTok{for}\NormalTok{ i }\KeywordTok{in} \BuiltInTok{range}\NormalTok{(}\BuiltInTok{len}\NormalTok{(nums)):}
            \ControlFlowTok{if}\NormalTok{ used[i]: }\ControlFlowTok{continue}
\NormalTok{            used[i] }\OperatorTok{=} \VariableTok{True}
\NormalTok{            path.append(nums[i])}
\NormalTok{            backtrack(path, used)}
\NormalTok{            path.pop()}
\NormalTok{            used[i] }\OperatorTok{=} \VariableTok{False}
\NormalTok{    backtrack([], [}\VariableTok{False}\NormalTok{] }\OperatorTok{*} \BuiltInTok{len}\NormalTok{(nums))}
    \ControlFlowTok{return}\NormalTok{ result}

\CommentTok{\# N{-}Queens:}
\KeywordTok{def}\NormalTok{ solve\_n\_queens(n):}
\NormalTok{    result }\OperatorTok{=}\NormalTok{ []}
\NormalTok{    cols }\OperatorTok{=} \BuiltInTok{set}\NormalTok{()}\OperatorTok{;}\NormalTok{ diag1 }\OperatorTok{=} \BuiltInTok{set}\NormalTok{()}\OperatorTok{;}\NormalTok{ diag2 }\OperatorTok{=} \BuiltInTok{set}\NormalTok{()}
\NormalTok{    board }\OperatorTok{=}\NormalTok{ [[}\StringTok{\textquotesingle{}.\textquotesingle{}} \ControlFlowTok{for}\NormalTok{ \_ }\KeywordTok{in} \BuiltInTok{range}\NormalTok{(n)] }\ControlFlowTok{for}\NormalTok{ \_ }\KeywordTok{in} \BuiltInTok{range}\NormalTok{(n)]}
    \KeywordTok{def}\NormalTok{ backtrack(row):}
        \ControlFlowTok{if}\NormalTok{ row }\OperatorTok{==}\NormalTok{ n:}
\NormalTok{            result.append([}\StringTok{\textquotesingle{}\textquotesingle{}}\NormalTok{.join(r) }\ControlFlowTok{for}\NormalTok{ r }\KeywordTok{in}\NormalTok{ board])}
            \ControlFlowTok{return}
        \ControlFlowTok{for}\NormalTok{ col }\KeywordTok{in} \BuiltInTok{range}\NormalTok{(n):}
            \ControlFlowTok{if}\NormalTok{ col }\KeywordTok{in}\NormalTok{ cols }\KeywordTok{or}\NormalTok{ (row}\OperatorTok{{-}}\NormalTok{col) }\KeywordTok{in}\NormalTok{ diag1 }\KeywordTok{or}\NormalTok{ (row}\OperatorTok{+}\NormalTok{col) }\KeywordTok{in}\NormalTok{ diag2:}
                \ControlFlowTok{continue}
\NormalTok{            cols.add(col)}\OperatorTok{;}\NormalTok{ diag1.add(row}\OperatorTok{{-}}\NormalTok{col)}\OperatorTok{;}\NormalTok{ diag2.add(row}\OperatorTok{+}\NormalTok{col)}
\NormalTok{            board[row][col] }\OperatorTok{=} \StringTok{\textquotesingle{}Q\textquotesingle{}}
\NormalTok{            backtrack(row }\OperatorTok{+} \DecValTok{1}\NormalTok{)}
\NormalTok{            cols.remove(col)}\OperatorTok{;}\NormalTok{ diag1.remove(row}\OperatorTok{{-}}\NormalTok{col)}\OperatorTok{;}\NormalTok{ diag2.remove(row}\OperatorTok{+}\NormalTok{col)}
\NormalTok{            board[row][col] }\OperatorTok{=} \StringTok{\textquotesingle{}.\textquotesingle{}}
\NormalTok{    backtrack(}\DecValTok{0}\NormalTok{)}
    \ControlFlowTok{return}\NormalTok{ result}
\end{Highlighting}
\end{Shaded}

\textbf{Pruning strategies:}

\begin{itemize}
\tightlist
\item
  \textbf{Constraint propagation:} Skip branches that violate constraints immediately
\item
  \textbf{Bound pruning:} Skip branches that cannot improve on best known solution
\item
  \textbf{Symmetry breaking:} Avoid exploring symmetrically equivalent states
\item
  \textbf{Early termination:} Return as soon as one solution found (if only need one)
\end{itemize}

\textbf{Backtracking complexity:} Exponential in worst case (O(n!), O(2\(^n\))), but pruning can dramatically reduce actual runtime.

\textbf{Backtracking is NOT DP:} Backtracking explores all possibilities with pruning; DP stores subproblem results to avoid recomputation. Backtracking problems don\textquotesingle t have overlapping subproblems.

\subsection{Bit Manipulation}\label{bit-manipulation}

Computers operate on binary; bit tricks exploit this for O(1) operations.

\subsubsection{Core Operations}\label{core-operations}

\begin{longtable}[]{@{}
  >{\raggedright\arraybackslash}p{(\columnwidth - 4\tabcolsep) * \real{0.2619}}
  >{\raggedright\arraybackslash}p{(\columnwidth - 4\tabcolsep) * \real{0.4048}}
  >{\raggedright\arraybackslash}p{(\columnwidth - 4\tabcolsep) * \real{0.3333}}@{}}
\toprule\noalign{}
\rowcolor{acc!16}
\begin{minipage}[b]{\linewidth}\raggedright
Operation
\end{minipage} & \begin{minipage}[b]{\linewidth}\raggedright
Expression
\end{minipage} & \begin{minipage}[b]{\linewidth}\raggedright
Effect
\end{minipage} \\
\midrule\noalign{}
\endhead
\bottomrule\noalign{}
\endlastfoot
Set bit k & `n & = (1 \textless\textless{} k)` \\
\rowcolor{zebra}
Clear bit k & \texttt{n\ \&=\ \textasciitilde{}(1\ \textless{}\textless{}\ k)} & Force bit k to 0 \\
Toggle bit k & \texttt{n\ \^{}=\ (1\ \textless{}\textless{}\ k)} & Flip bit k \\
\rowcolor{zebra}
Check bit k & \texttt{(n\ \textgreater{}\textgreater{}\ k)\ \&\ 1} & Get value of bit k \\
Clear lowest set bit & \texttt{n\ \&\ (n\ -\ 1)} & Removes rightmost 1 bit \\
\rowcolor{zebra}
Isolate lowest set bit & \texttt{n\ \&\ (-n)} & Gets rightmost 1 bit only \\
Check power of 2 & \texttt{n\ \textgreater{}\ 0\ and\ (n\ \&\ (n-1))\ ==\ 0} & True iff exactly one bit set \\
\rowcolor{zebra}
Count set bits & \texttt{bin(n).count(\textquotesingle{}1\textquotesingle{})} or \texttt{n.bit\_count()} & Hamming weight \\
\end{longtable}

\subsubsection{Key XOR Tricks}\label{key-xor-tricks}

\setcodelang{Python}

\begin{Shaded}
\begin{Highlighting}[]
\CommentTok{\# XOR properties:}
\CommentTok{\# a \^{} a = 0 (anything XOR itself = 0)}
\CommentTok{\# a \^{} 0 = a (anything XOR 0 = itself)}
\CommentTok{\# XOR is commutative and associative}

\CommentTok{\# Find single number in array where all others appear twice:}
\KeywordTok{def}\NormalTok{ single\_number(nums):}
\NormalTok{    result }\OperatorTok{=} \DecValTok{0}
    \ControlFlowTok{for}\NormalTok{ x }\KeywordTok{in}\NormalTok{ nums:}
\NormalTok{        result }\OperatorTok{\^{}=}\NormalTok{ x   }\CommentTok{\# pairs cancel out}
    \ControlFlowTok{return}\NormalTok{ result}

\CommentTok{\# Find two unique numbers when all others appear twice:}
\KeywordTok{def}\NormalTok{ single\_number\_two(nums):}
\NormalTok{    xor }\OperatorTok{=} \DecValTok{0}
    \ControlFlowTok{for}\NormalTok{ x }\KeywordTok{in}\NormalTok{ nums: xor }\OperatorTok{\^{}=}\NormalTok{ x   }\CommentTok{\# xor of the two unique numbers}
    \CommentTok{\# Find any set bit in xor (they differ in this bit)}
\NormalTok{    diff\_bit }\OperatorTok{=}\NormalTok{ xor }\OperatorTok{\&}\NormalTok{ (}\OperatorTok{{-}}\NormalTok{xor)}
\NormalTok{    a }\OperatorTok{=} \DecValTok{0}
    \ControlFlowTok{for}\NormalTok{ x }\KeywordTok{in}\NormalTok{ nums:}
        \ControlFlowTok{if}\NormalTok{ x }\OperatorTok{\&}\NormalTok{ diff\_bit:}
\NormalTok{            a }\OperatorTok{\^{}=}\NormalTok{ x   }\CommentTok{\# group by this bit}
    \ControlFlowTok{return}\NormalTok{ a, xor }\OperatorTok{\^{}}\NormalTok{ a}

\CommentTok{\# Swap without temp:}
\NormalTok{a }\OperatorTok{\^{}=}\NormalTok{ b}\OperatorTok{;}\NormalTok{ b }\OperatorTok{\^{}=}\NormalTok{ a}\OperatorTok{;}\NormalTok{ a }\OperatorTok{\^{}=}\NormalTok{ b}

\CommentTok{\# Check if two integers have opposite signs:}
\KeywordTok{def}\NormalTok{ opposite\_signs(a, b):}
    \ControlFlowTok{return}\NormalTok{ (a }\OperatorTok{\^{}}\NormalTok{ b) }\OperatorTok{\textless{}} \DecValTok{0}   \CommentTok{\# MSB will be 1 if signs differ}
\end{Highlighting}
\end{Shaded}

\subsubsection{Bitmask DP}\label{bitmask-dp}

\setcodelang{Python}

\begin{Shaded}
\begin{Highlighting}[]
\CommentTok{\# Enumerate all subsets of a bitmask:}
\NormalTok{mask }\OperatorTok{=} \BaseNTok{0b1011}   \CommentTok{\# some bitmask}
\NormalTok{sub }\OperatorTok{=}\NormalTok{ mask}
\ControlFlowTok{while}\NormalTok{ sub }\OperatorTok{\textgreater{}} \DecValTok{0}\NormalTok{:}
\NormalTok{    process(sub)   }\CommentTok{\# sub is a subset of mask}
\NormalTok{    sub }\OperatorTok{=}\NormalTok{ (sub }\OperatorTok{{-}} \DecValTok{1}\NormalTok{) }\OperatorTok{\&}\NormalTok{ mask   }\CommentTok{\# get next subset}

\CommentTok{\# Check if bit i is set in mask:}
\ControlFlowTok{if}\NormalTok{ mask }\OperatorTok{\&}\NormalTok{ (}\DecValTok{1} \OperatorTok{\textless{}\textless{}}\NormalTok{ i):}
    \ControlFlowTok{pass}

\CommentTok{\# Count set bits (Brian Kernighan\textquotesingle{}s algorithm):}
\KeywordTok{def}\NormalTok{ count\_bits(n):}
\NormalTok{    count }\OperatorTok{=} \DecValTok{0}
    \ControlFlowTok{while}\NormalTok{ n:}
\NormalTok{        n }\OperatorTok{\&=}\NormalTok{ (n }\OperatorTok{{-}} \DecValTok{1}\NormalTok{)   }\CommentTok{\# remove lowest set bit}
\NormalTok{        count }\OperatorTok{+=} \DecValTok{1}
    \ControlFlowTok{return}\NormalTok{ count}

\CommentTok{\# Generate all 2\^{}n subsets:}
\NormalTok{n }\OperatorTok{=} \DecValTok{4}
\ControlFlowTok{for}\NormalTok{ mask }\KeywordTok{in} \BuiltInTok{range}\NormalTok{(}\DecValTok{1} \OperatorTok{\textless{}\textless{}}\NormalTok{ n):}
\NormalTok{    subset }\OperatorTok{=}\NormalTok{ [i }\ControlFlowTok{for}\NormalTok{ i }\KeywordTok{in} \BuiltInTok{range}\NormalTok{(n) }\ControlFlowTok{if}\NormalTok{ mask }\OperatorTok{\&}\NormalTok{ (}\DecValTok{1} \OperatorTok{\textless{}\textless{}}\NormalTok{ i)]}
\end{Highlighting}
\end{Shaded}

\subsubsection{Bit Tricks Summary}\label{bit-tricks-summary}

\setcodelang{Python}

\begin{Shaded}
\begin{Highlighting}[]
\CommentTok{\# Multiply/divide by power of 2:}
\NormalTok{x }\OperatorTok{\textless{}\textless{}}\NormalTok{ k  }\OperatorTok{==}\NormalTok{ x }\OperatorTok{*} \DecValTok{2}\OperatorTok{\^{}}\NormalTok{k}
\NormalTok{x }\OperatorTok{\textgreater{}\textgreater{}}\NormalTok{ k  }\OperatorTok{==}\NormalTok{ x }\OperatorTok{//} \DecValTok{2}\OperatorTok{\^{}}\NormalTok{k}

\CommentTok{\# Modulo by power of 2:}
\NormalTok{x }\OperatorTok{\&}\NormalTok{ (m }\OperatorTok{{-}} \DecValTok{1}\NormalTok{) }\OperatorTok{==}\NormalTok{ x }\OperatorTok{\%}\NormalTok{ m   }\CommentTok{\# only works when m is power of 2}

\CommentTok{\# Average without overflow:}
\NormalTok{avg }\OperatorTok{=}\NormalTok{ (a }\OperatorTok{\&}\NormalTok{ b) }\OperatorTok{+}\NormalTok{ ((a }\OperatorTok{\^{}}\NormalTok{ b) }\OperatorTok{\textgreater{}\textgreater{}} \DecValTok{1}\NormalTok{)}

\CommentTok{\# Absolute value (no branching):}
\NormalTok{mask }\OperatorTok{=}\NormalTok{ x }\OperatorTok{\textgreater{}\textgreater{}} \DecValTok{31}   \CommentTok{\# {-}1 if negative, 0 if positive}
\NormalTok{abs\_x }\OperatorTok{=}\NormalTok{ (x }\OperatorTok{+}\NormalTok{ mask) }\OperatorTok{\^{}}\NormalTok{ mask}

\CommentTok{\# Next power of 2:}
\KeywordTok{def}\NormalTok{ next\_power\_of\_2(n):}
\NormalTok{    n }\OperatorTok{{-}=} \DecValTok{1}
\NormalTok{    n }\OperatorTok{|=}\NormalTok{ n }\OperatorTok{\textgreater{}\textgreater{}} \DecValTok{1}\OperatorTok{;}\NormalTok{ n }\OperatorTok{|=}\NormalTok{ n }\OperatorTok{\textgreater{}\textgreater{}} \DecValTok{2}\OperatorTok{;}\NormalTok{ n }\OperatorTok{|=}\NormalTok{ n }\OperatorTok{\textgreater{}\textgreater{}} \DecValTok{4}
\NormalTok{    n }\OperatorTok{|=}\NormalTok{ n }\OperatorTok{\textgreater{}\textgreater{}} \DecValTok{8}\OperatorTok{;}\NormalTok{ n }\OperatorTok{|=}\NormalTok{ n }\OperatorTok{\textgreater{}\textgreater{}} \DecValTok{16}
    \ControlFlowTok{return}\NormalTok{ n }\OperatorTok{+} \DecValTok{1}
\end{Highlighting}
\end{Shaded}

\subsection{Number Theory / Math}\label{number-theory--math}

\subsubsection{GCD and LCM}\label{gcd-and-lcm}

\setcodelang{Python}

\begin{Shaded}
\begin{Highlighting}[]
\ImportTok{import}\NormalTok{ math}

\CommentTok{\# GCD using Euclidean algorithm: gcd(a, b) = gcd(b, a \% b)}
\KeywordTok{def}\NormalTok{ gcd(a, b):}
    \ControlFlowTok{while}\NormalTok{ b:}
\NormalTok{        a, b }\OperatorTok{=}\NormalTok{ b, a }\OperatorTok{\%}\NormalTok{ b}
    \ControlFlowTok{return}\NormalTok{ a}

\CommentTok{\# LCM: lcm(a, b) = a * b / gcd(a, b)}
\KeywordTok{def}\NormalTok{ lcm(a, b):}
    \ControlFlowTok{return}\NormalTok{ a }\OperatorTok{*}\NormalTok{ b }\OperatorTok{//}\NormalTok{ gcd(a, b)}

\CommentTok{\# Python 3.9+:}
\NormalTok{math.gcd(a, b)}
\NormalTok{math.lcm(a, b)}

\CommentTok{\# GCD complexity: O(log(min(a, b))) — Fibonacci numbers are worst case}
\end{Highlighting}
\end{Shaded}

\textbf{Extended Euclidean Algorithm:} Finds x, y such that a\(\cdot\)x + b\(\cdot\)y = gcd(a, b):

\setcodelang{Python}

\begin{Shaded}
\begin{Highlighting}[]
\KeywordTok{def}\NormalTok{ extended\_gcd(a, b):}
    \ControlFlowTok{if}\NormalTok{ b }\OperatorTok{==} \DecValTok{0}\NormalTok{:}
        \ControlFlowTok{return}\NormalTok{ a, }\DecValTok{1}\NormalTok{, }\DecValTok{0}
\NormalTok{    g, x, y }\OperatorTok{=}\NormalTok{ extended\_gcd(b, a }\OperatorTok{\%}\NormalTok{ b)}
    \ControlFlowTok{return}\NormalTok{ g, y, x }\OperatorTok{{-}}\NormalTok{ (a }\OperatorTok{//}\NormalTok{ b) }\OperatorTok{*}\NormalTok{ y}
\end{Highlighting}
\end{Shaded}

\subsubsection{Sieve of Eratosthenes}\label{sieve-of-eratosthenes}

Find all primes up to n in O(n log log n):

\setcodelang{Python}

\begin{Shaded}
\begin{Highlighting}[]
\KeywordTok{def}\NormalTok{ sieve(n):}
\NormalTok{    is\_prime }\OperatorTok{=}\NormalTok{ [}\VariableTok{True}\NormalTok{] }\OperatorTok{*}\NormalTok{ (n }\OperatorTok{+} \DecValTok{1}\NormalTok{)}
\NormalTok{    is\_prime[}\DecValTok{0}\NormalTok{] }\OperatorTok{=}\NormalTok{ is\_prime[}\DecValTok{1}\NormalTok{] }\OperatorTok{=} \VariableTok{False}
    \ControlFlowTok{for}\NormalTok{ i }\KeywordTok{in} \BuiltInTok{range}\NormalTok{(}\DecValTok{2}\NormalTok{, }\BuiltInTok{int}\NormalTok{(n}\OperatorTok{**}\FloatTok{0.5}\NormalTok{) }\OperatorTok{+} \DecValTok{1}\NormalTok{):}
        \ControlFlowTok{if}\NormalTok{ is\_prime[i]:}
            \ControlFlowTok{for}\NormalTok{ j }\KeywordTok{in} \BuiltInTok{range}\NormalTok{(i}\OperatorTok{*}\NormalTok{i, n}\OperatorTok{+}\DecValTok{1}\NormalTok{, i):   }\CommentTok{\# start from i² (smaller multiples already marked)}
\NormalTok{                is\_prime[j] }\OperatorTok{=} \VariableTok{False}
    \ControlFlowTok{return}\NormalTok{ [i }\ControlFlowTok{for}\NormalTok{ i }\KeywordTok{in} \BuiltInTok{range}\NormalTok{(n}\OperatorTok{+}\DecValTok{1}\NormalTok{) }\ControlFlowTok{if}\NormalTok{ is\_prime[i]]}

\CommentTok{\# Segmented sieve for large ranges:}
\CommentTok{\# Prime factorization using sieve (smallest prime factor):}
\KeywordTok{def}\NormalTok{ smallest\_prime\_factor(n):}
\NormalTok{    spf }\OperatorTok{=} \BuiltInTok{list}\NormalTok{(}\BuiltInTok{range}\NormalTok{(n }\OperatorTok{+} \DecValTok{1}\NormalTok{))}
    \ControlFlowTok{for}\NormalTok{ i }\KeywordTok{in} \BuiltInTok{range}\NormalTok{(}\DecValTok{2}\NormalTok{, }\BuiltInTok{int}\NormalTok{(n}\OperatorTok{**}\FloatTok{0.5}\NormalTok{) }\OperatorTok{+} \DecValTok{1}\NormalTok{):}
        \ControlFlowTok{if}\NormalTok{ spf[i] }\OperatorTok{==}\NormalTok{ i:   }\CommentTok{\# i is prime}
            \ControlFlowTok{for}\NormalTok{ j }\KeywordTok{in} \BuiltInTok{range}\NormalTok{(i}\OperatorTok{*}\NormalTok{i, n}\OperatorTok{+}\DecValTok{1}\NormalTok{, i):}
                \ControlFlowTok{if}\NormalTok{ spf[j] }\OperatorTok{==}\NormalTok{ j:}
\NormalTok{                    spf[j] }\OperatorTok{=}\NormalTok{ i}
    \ControlFlowTok{return}\NormalTok{ spf}
\end{Highlighting}
\end{Shaded}

\subsubsection{Modular Arithmetic}\label{modular-arithmetic}

For large number problems where answer is \texttt{mod\ 10\^{}9\ +\ 7}:

\setcodelang{Python}

\begin{Shaded}
\begin{Highlighting}[]
\NormalTok{MOD }\OperatorTok{=} \DecValTok{10}\OperatorTok{**}\DecValTok{9} \OperatorTok{+} \DecValTok{7}

\CommentTok{\# Basic operations:}
\NormalTok{(a }\OperatorTok{+}\NormalTok{ b) }\OperatorTok{\%}\NormalTok{ MOD}
\NormalTok{(a }\OperatorTok{{-}}\NormalTok{ b }\OperatorTok{+}\NormalTok{ MOD) }\OperatorTok{\%}\NormalTok{ MOD   }\CommentTok{\# add MOD before \% to avoid negative}
\NormalTok{(a }\OperatorTok{*}\NormalTok{ b) }\OperatorTok{\%}\NormalTok{ MOD}
\NormalTok{(a }\OperatorTok{**}\NormalTok{ b) }\OperatorTok{\%}\NormalTok{ MOD }\OperatorTok{==} \BuiltInTok{pow}\NormalTok{(a, b, MOD)   }\CommentTok{\# Python\textquotesingle{}s built{-}in is efficient}

\CommentTok{\# Modular inverse (when MOD is prime, use Fermat\textquotesingle{}s little theorem):}
\CommentTok{\# a\^{}(MOD{-}1) ≡ 1 (mod MOD) when MOD is prime}
\CommentTok{\# Therefore a\^{}({-}1) ≡ a\^{}(MOD{-}2) (mod MOD)}
\KeywordTok{def}\NormalTok{ mod\_inverse(a, mod):}
    \ControlFlowTok{return} \BuiltInTok{pow}\NormalTok{(a, mod }\OperatorTok{{-}} \DecValTok{2}\NormalTok{, mod)   }\CommentTok{\# O(log mod)}

\CommentTok{\# Modular inverse using extended GCD (works for any mod, gcd(a,mod)=1):}
\KeywordTok{def}\NormalTok{ mod\_inverse\_gcd(a, mod):}
\NormalTok{    g, x, \_ }\OperatorTok{=}\NormalTok{ extended\_gcd(a, mod)}
    \ControlFlowTok{if}\NormalTok{ g }\OperatorTok{!=} \DecValTok{1}\NormalTok{: }\ControlFlowTok{raise} \PreprocessorTok{ValueError}\NormalTok{(}\StringTok{"Inverse doesn\textquotesingle{}t exist"}\NormalTok{)}
    \ControlFlowTok{return}\NormalTok{ x }\OperatorTok{\%}\NormalTok{ mod}

\CommentTok{\# Modular combinations:}
\KeywordTok{def}\NormalTok{ mod\_choose(n, k, mod}\OperatorTok{=}\DecValTok{10}\OperatorTok{**}\DecValTok{9}\OperatorTok{+}\DecValTok{7}\NormalTok{):}
    \ControlFlowTok{if}\NormalTok{ k }\OperatorTok{\textgreater{}}\NormalTok{ n: }\ControlFlowTok{return} \DecValTok{0}
    \CommentTok{\# Precompute factorials}
\NormalTok{    fact }\OperatorTok{=}\NormalTok{ [}\DecValTok{1}\NormalTok{] }\OperatorTok{*}\NormalTok{ (n }\OperatorTok{+} \DecValTok{1}\NormalTok{)}
    \ControlFlowTok{for}\NormalTok{ i }\KeywordTok{in} \BuiltInTok{range}\NormalTok{(}\DecValTok{1}\NormalTok{, n}\OperatorTok{+}\DecValTok{1}\NormalTok{):}
\NormalTok{        fact[i] }\OperatorTok{=}\NormalTok{ fact[i}\OperatorTok{{-}}\DecValTok{1}\NormalTok{] }\OperatorTok{*}\NormalTok{ i }\OperatorTok{\%}\NormalTok{ mod}
\NormalTok{    inv\_fact }\OperatorTok{=}\NormalTok{ [}\DecValTok{1}\NormalTok{] }\OperatorTok{*}\NormalTok{ (n }\OperatorTok{+} \DecValTok{1}\NormalTok{)}
\NormalTok{    inv\_fact[n] }\OperatorTok{=} \BuiltInTok{pow}\NormalTok{(fact[n], mod }\OperatorTok{{-}} \DecValTok{2}\NormalTok{, mod)}
    \ControlFlowTok{for}\NormalTok{ i }\KeywordTok{in} \BuiltInTok{range}\NormalTok{(n}\OperatorTok{{-}}\DecValTok{1}\NormalTok{, }\OperatorTok{{-}}\DecValTok{1}\NormalTok{, }\OperatorTok{{-}}\DecValTok{1}\NormalTok{):}
\NormalTok{        inv\_fact[i] }\OperatorTok{=}\NormalTok{ inv\_fact[i}\OperatorTok{+}\DecValTok{1}\NormalTok{] }\OperatorTok{*}\NormalTok{ (i}\OperatorTok{+}\DecValTok{1}\NormalTok{) }\OperatorTok{\%}\NormalTok{ mod}
    \ControlFlowTok{return}\NormalTok{ fact[n] }\OperatorTok{*}\NormalTok{ inv\_fact[k] }\OperatorTok{\%}\NormalTok{ mod }\OperatorTok{*}\NormalTok{ inv\_fact[n}\OperatorTok{{-}}\NormalTok{k] }\OperatorTok{\%}\NormalTok{ mod}
\end{Highlighting}
\end{Shaded}

\subsubsection{Fast Exponentiation (Binary Exponentiation)}\label{fast-exponentiation-binary-exponentiation}

Compute a\^{}b in O(log b):

\setcodelang{Python}

\begin{Shaded}
\begin{Highlighting}[]
\KeywordTok{def}\NormalTok{ fast\_pow(base, exp, mod}\OperatorTok{=}\VariableTok{None}\NormalTok{):}
\NormalTok{    result }\OperatorTok{=} \DecValTok{1}
    \ControlFlowTok{while}\NormalTok{ exp }\OperatorTok{\textgreater{}} \DecValTok{0}\NormalTok{:}
        \ControlFlowTok{if}\NormalTok{ exp }\OperatorTok{\&} \DecValTok{1}\NormalTok{:               }\CommentTok{\# if current bit is 1}
\NormalTok{            result }\OperatorTok{=}\NormalTok{ result }\OperatorTok{*}\NormalTok{ base}
            \ControlFlowTok{if}\NormalTok{ mod: result }\OperatorTok{\%=}\NormalTok{ mod}
\NormalTok{        base }\OperatorTok{=}\NormalTok{ base }\OperatorTok{*}\NormalTok{ base}
        \ControlFlowTok{if}\NormalTok{ mod: base }\OperatorTok{\%=}\NormalTok{ mod}
\NormalTok{        exp }\OperatorTok{\textgreater{}\textgreater{}=} \DecValTok{1}                 \CommentTok{\# shift right (divide by 2)}
    \ControlFlowTok{return}\NormalTok{ result}

\CommentTok{\# Python\textquotesingle{}s built{-}in pow(a, b, mod) uses this algorithm}
\end{Highlighting}
\end{Shaded}

\textbf{Key idea:} a\^{}13 = a\^{}(1101 in binary) = a\^{}8 \(\cdot\) a\^{}4 \(\cdot\) a\^{}1. Square the base repeatedly, multiply result when bit is 1.

\subsubsection{Prime Factorization}\label{prime-factorization}

\setcodelang{Python}

\begin{Shaded}
\begin{Highlighting}[]
\KeywordTok{def}\NormalTok{ prime\_factors(n):}
\NormalTok{    factors }\OperatorTok{=}\NormalTok{ \{\}}
\NormalTok{    d }\OperatorTok{=} \DecValTok{2}
    \ControlFlowTok{while}\NormalTok{ d }\OperatorTok{*}\NormalTok{ d }\OperatorTok{\textless{}=}\NormalTok{ n:}
        \ControlFlowTok{while}\NormalTok{ n }\OperatorTok{\%}\NormalTok{ d }\OperatorTok{==} \DecValTok{0}\NormalTok{:}
\NormalTok{            factors[d] }\OperatorTok{=}\NormalTok{ factors.get(d, }\DecValTok{0}\NormalTok{) }\OperatorTok{+} \DecValTok{1}
\NormalTok{            n }\OperatorTok{//=}\NormalTok{ d}
\NormalTok{        d }\OperatorTok{+=} \DecValTok{1}
    \ControlFlowTok{if}\NormalTok{ n }\OperatorTok{\textgreater{}} \DecValTok{1}\NormalTok{:}
\NormalTok{        factors[n] }\OperatorTok{=}\NormalTok{ factors.get(n, }\DecValTok{0}\NormalTok{) }\OperatorTok{+} \DecValTok{1}
    \ControlFlowTok{return}\NormalTok{ factors}
\CommentTok{\# Complexity: O(√n)}
\end{Highlighting}
\end{Shaded}

\subsubsection{Useful Math Facts}\label{useful-math-facts}

\begin{itemize}
\tightlist
\item
  Number of divisors of n: O(n\^{}(1/3)) on average
\item
  Sum of 1 to n: n(n+1)/2
\item
  Sum of squares 1 to n: n(n+1)(2n+1)/6
\item
  Catalan numbers: C(2n,n)/(n+1) --- counts BSTs, valid parentheses, polygon triangulations
\item
  \textbf{Pigeonhole principle:} n+1 items in n containers \(\rightarrow\) at least one container has \(\geq\) 2 items
\item
  \textbf{Stars and bars:} Distribute n identical items into k bins: C(n+k-1, k-1)
\end{itemize}

\section{Algorithm Selection Cheat Table}\label{algorithm-selection-cheat-table}

\begin{longtable}[]{@{}
  >{\raggedright\arraybackslash}p{(\columnwidth - 6\tabcolsep) * \real{0.2991}}
  >{\raggedright\arraybackslash}p{(\columnwidth - 6\tabcolsep) * \real{0.2564}}
  >{\raggedright\arraybackslash}p{(\columnwidth - 6\tabcolsep) * \real{0.2222}}
  >{\raggedright\arraybackslash}p{(\columnwidth - 6\tabcolsep) * \real{0.2222}}@{}}
\toprule\noalign{}
\rowcolor{acc!16}
\begin{minipage}[b]{\linewidth}\raggedright
Problem Type
\end{minipage} & \begin{minipage}[b]{\linewidth}\raggedright
Key Signal
\end{minipage} & \begin{minipage}[b]{\linewidth}\raggedright
Recommended Algorithm
\end{minipage} & \begin{minipage}[b]{\linewidth}\raggedright
Complexity
\end{minipage} \\
\midrule\noalign{}
\endhead
\bottomrule\noalign{}
\endlastfoot
Shortest path, unweighted & BFS in grid/graph & BFS & O(V+E) \\
\rowcolor{zebra}
Shortest path, non-negative weights & Edge weights, single source & Dijkstra & O((V+E)log V) \\
Shortest path, negative weights & Negative edges allowed & Bellman-Ford & O(VE) \\
\rowcolor{zebra}
All-pairs shortest path & Need dist between all pairs & Floyd-Warshall & O(V\(^3\)) \\
Minimum spanning tree & Connect all nodes cheaply & Kruskal or Prim & O(E log E) \\
\rowcolor{zebra}
Sort, general & n \(\leq\) 10\^{}6 & Quick sort / Merge sort & O(n log n) \\
Sort, small integer range & Keys in {[}0, k{]} & Counting sort / Radix sort & O(n+k) \\
\rowcolor{zebra}
Cycle detection, undirected & Find cycle in graph & Union-Find or DFS & O(\(\alpha\)(n)) or O(V+E) \\
Cycle detection, directed & Find cycle in digraph & DFS (3-color) & O(V+E) \\
\rowcolor{zebra}
Topological order & DAG dependency order & Kahn\textquotesingle s or DFS topo & O(V+E) \\
Connected components & Group nodes & DFS/BFS + visited set & O(V+E) \\
\rowcolor{zebra}
Optimal value with choices & Overlapping subproblems & Dynamic Programming & Varies \\
Make decisions greedily & Exchange argument holds & Greedy & O(n log n) or O(n) \\
\rowcolor{zebra}
All subsets/permutations & Enumerate possibilities & Backtracking & O(2\^{}n) or O(n!) \\
Find element in sorted array & Array is sorted & Binary search & O(log n) \\
\rowcolor{zebra}
Find minimum X satisfying condition & Monotone predicate on integers & Binary search on answer & O(log n \(\cdot\) check) \\
Frequency counting & Count occurrences & Hash map / Counting array & O(n) \\
\rowcolor{zebra}
Range minimum/maximum & Static range queries & Sparse table & O(n log n) pre, O(1) query \\
Range sum queries & Prefix sums or updates & Prefix sum / Fenwick tree & O(n) or O(n log n) \\
\rowcolor{zebra}
Substring search & Pattern in text & KMP / Rabin-Karp & O(n+m) \\
Disjoint set operations & Dynamic connectivity & Union-Find & O(\(\alpha\)(n)) amortized \\
\rowcolor{zebra}
Large integers, divisibility & Modular arithmetic & Fermat\textquotesingle s theorem, sieve & O(log n) \\
Find unique / pair in XOR & XOR properties & Bit manipulation & O(n) \\
\rowcolor{zebra}
Multiple choices at each step & Combinatorial search & Backtracking + pruning & Exponential \\
Maximize non-overlapping intervals & Interval scheduling & Greedy (sort by end time) & O(n log n) \\
\end{longtable}

\section{Real-Life Analogies}\label{real-life-analogies}

\begin{longtable}[]{@{}
  >{\raggedright\arraybackslash}p{(\columnwidth - 2\tabcolsep) * \real{0.3239}}
  >{\raggedright\arraybackslash}p{(\columnwidth - 2\tabcolsep) * \real{0.6761}}@{}}
\toprule\noalign{}
\rowcolor{acc!16}
\begin{minipage}[b]{\linewidth}\raggedright
Algorithm / Concept
\end{minipage} & \begin{minipage}[b]{\linewidth}\raggedright
Real-Life Analogy
\end{minipage} \\
\midrule\noalign{}
\endhead
\bottomrule\noalign{}
\endlastfoot
BFS & Ripple in a pond --- spreads layer by layer from center \\
\rowcolor{zebra}
DFS & Solving a maze by always turning left until you hit a wall, then backtracking \\
Dijkstra & Google Maps routing --- always expand the nearest unvisited city \\
\rowcolor{zebra}
Bellman-Ford & Send rumors through a network; after n-1 rounds, everyone knows the shortest path \\
Dynamic Programming & Filling out a tax form --- use last year\textquotesingle s answers to compute this year\textquotesingle s \\
\rowcolor{zebra}
Memoization & Sticky notes --- write down each sub-answer once so you never recalculate it \\
Greedy (interval scheduling) & Booking meeting rooms --- always pick the meeting that ends soonest \\
\rowcolor{zebra}
Merge Sort & Library sorting books by splitting into piles, sorting each pile, then merging \\
Quick Sort & Separating audience by height relative to a pivot person, then recursing \\
\rowcolor{zebra}
Binary Search & Guessing a number --- always guess the middle of the remaining range \\
Topological Sort & Course prerequisite ordering --- take prerequisites before the main course \\
\rowcolor{zebra}
Union-Find & Friend groups --- if A and B are friends, merge their groups; check if same group \\
Backtracking & Trying all keys on a keyring --- discard a key as soon as it doesn\textquotesingle t fit \\
\rowcolor{zebra}
Sieve of Eratosthenes & Crossing out multiples on a number grid to find primes \\
Amortized Analysis & Buying coffee with a punch card --- most coffees are cheap, the 10th is free \\
\rowcolor{zebra}
Big-O & Speed limit signs --- tells you the worst that can happen, not typical speed \\
MST (Kruskal) & Building cheapest road network by adding roads from cheapest to most expensive \\
\rowcolor{zebra}
Hash Table & Lockers with assigned numbers --- direct access by computing the locker number \\
Heap & Hospital triage --- the most critical patient always treated next \\
\end{longtable}

\section{Common Interview Questions}\label{common-interview-questions}

\subsection{Complexity Analysis}\label{complexity-analysis-1}

\begin{itemize}
\tightlist
\item
  What is the time complexity of Python\textquotesingle s \texttt{list.sort()}? (O(n log n) --- Tim Sort)
\item
  What is the amortized complexity of appending to a Python list? (O(1))
\item
  Given a recursive function, derive its time complexity using Master Theorem
\item
  What is the space complexity of BFS vs DFS? (BFS: O(V) for queue; DFS: O(V) for stack/recursion depth)
\end{itemize}

\subsection{Sorting}\label{sorting-1}

\begin{itemize}
\tightlist
\item
  Sort an array of 0s, 1s, and 2s without using built-in sort (Dutch National Flag --- O(n))
\item
  Find the k-th largest element in an array (QuickSelect --- O(n) avg, or heap --- O(n log k))
\item
  Check if an array is sorted in O(n)
\item
  Merge two sorted arrays in-place
\end{itemize}

\subsection{Binary Search}\label{binary-search-1}

\begin{itemize}
\tightlist
\item
  Find first and last position of element in sorted array (LeetCode 34)
\item
  Search in rotated sorted array (LeetCode 33)
\item
  Find minimum in rotated sorted array (LeetCode 153)
\item
  Koko eating bananas / minimum capacity ship (binary search on answer space)
\item
  Find peak element (LeetCode 162 --- binary search on unsorted array!)
\end{itemize}

\subsection{Graphs}\label{graphs}

\begin{itemize}
\tightlist
\item
  Number of islands (DFS/BFS on grid --- LeetCode 200)
\item
  Course schedule --- detect cycle in directed graph (LeetCode 207)
\item
  Word ladder --- BFS shortest path (LeetCode 127)
\item
  Clone graph (DFS with memo)
\item
  Cheapest flights within k stops (modified Dijkstra/Bellman-Ford)
\item
  Reconstruct itinerary (Eulerian path via DFS)
\item
  Network delay time (Dijkstra --- LeetCode 743)
\item
  Find if path exists --- Union-Find
\item
  Minimum cost to connect all points --- Prim\textquotesingle s MST (LeetCode 1584)
\item
  Critical connections / bridges (Tarjan\textquotesingle s algorithm)
\end{itemize}

\subsection{Dynamic Programming}\label{dynamic-programming-1}

\begin{itemize}
\tightlist
\item
  Climbing stairs / coin change (basic 1D DP)
\item
  Longest common subsequence (classic 2D DP)
\item
  Knapsack variants (0/1, unbounded, multi-dimensional)
\item
  Longest increasing subsequence (O(n\(^2\)) DP or O(n log n) patience sort)
\item
  Edit distance / regular expression matching
\item
  Maximum subarray sum (Kadane\textquotesingle s algorithm --- O(n) DP)
\item
  Maximum product subarray
\item
  House robber / house robber II (circular)
\item
  Unique paths / minimum path sum in grid
\item
  Burst balloons (interval DP --- LeetCode 312)
\item
  Word break (1D DP + hash set)
\item
  Partition equal subset sum (subset sum DP)
\end{itemize}

\subsection{Backtracking}\label{backtracking}

\begin{itemize}
\tightlist
\item
  N-Queens (LeetCode 51)
\item
  Sudoku solver (LeetCode 37)
\item
  Generate all valid parentheses (LeetCode 22)
\item
  Subsets / permutations / combinations (LeetCode 78, 46, 77)
\item
  Word search in grid (LeetCode 79)
\end{itemize}

\subsection{Greedy}\label{greedy}

\begin{itemize}
\tightlist
\item
  Merge overlapping intervals (LeetCode 56)
\item
  Non-overlapping intervals (LeetCode 435)
\item
  Jump game I and II (LeetCode 55, 45)
\item
  Task scheduler (LeetCode 621)
\item
  Minimum number of arrows to burst balloons (LeetCode 452)
\end{itemize}

\section{Typical Mistakes Candidates Make}\label{typical-mistakes-candidates-make}

\subsection{Complexity Analysis Mistakes}\label{complexity-analysis-mistakes}

\begin{itemize}
\tightlist
\item
  \textbf{Forgetting recursion call stack space} --- O(n) space for O(n) deep recursion
\item
  \textbf{Counting operations instead of asymptotic complexity} --- writing "O(2n)" instead of "O(n)"
\item
  \textbf{Wrong space complexity for BFS} --- BFS can hold O(V) nodes in queue, not O(1)
\item
  \textbf{Not accounting for sorting} --- saying a solution is O(n) when it first sorts in O(n log n)
\item
  \textbf{Misapplying Master Theorem} --- forgetting to check regularity condition for Case 3
\end{itemize}

\subsection{Graph Algorithm Mistakes}\label{graph-algorithm-mistakes}

\begin{itemize}
\tightlist
\item
  \textbf{Using Dijkstra with negative edges} --- it will produce wrong results
\item
  \textbf{Not handling disconnected graphs} --- BFS/DFS from single source misses other components
\item
  \textbf{Off-by-one in Bellman-Ford} --- need exactly V-1 iterations, not V
\item
  \textbf{Forgetting to check for cycles in topological sort} --- not verifying all nodes are processed
\item
  \textbf{Confusing directed and undirected cycle detection} --- using undirected method on directed graph
\item
  \textbf{Not initializing all distances to infinity} --- leaving 0 causes wrong Dijkstra results
\end{itemize}

\subsection{Sorting Mistakes}\label{sorting-mistakes}

\begin{itemize}
\tightlist
\item
  \textbf{Using comparison sort when counting/radix would be O(n)}
\item
  \textbf{Forgetting stability requirements} --- using unstable sort when stability matters
\item
  \textbf{Not randomizing quicksort pivot} --- gets O(n\(^2\)) on sorted input in interviews where test cases include sorted arrays
\item
  \textbf{Forgetting in-place constraint} --- using merge sort when O(1) space needed
\end{itemize}

\subsection{Dynamic Programming Mistakes}\label{dynamic-programming-mistakes}

\begin{itemize}
\tightlist
\item
  \textbf{Defining state incorrectly} --- most DP bugs come from wrong state definition
\item
  \textbf{Wrong iteration order} --- 1D knapsack must iterate weights in reverse for 0/1; forward for unbounded
\item
  \textbf{Off-by-one in indices} --- dp{[}i{]} represents item i or first i items?
\item
  \textbf{Not thinking about base cases} --- wrong base cases cause cascade of wrong answers
\item
  \textbf{Confusing DP with backtracking} --- implementing exponential backtracking when DP is needed
\item
  \textbf{Forgetting to return correct cell} --- computing dp{[}n{]}{[}m{]} when answer is in dp{[}n-1{]}{[}m-1{]}
\end{itemize}

\subsection{Binary Search Mistakes}\label{binary-search-mistakes}

\begin{itemize}
\tightlist
\item
  \textbf{Integer overflow with mid = (lo + hi) / 2} --- use \texttt{lo\ +\ (hi\ -\ lo)\ //\ 2}
\item
  \textbf{Infinite loop} --- wrong loop condition or mid not shrinking the range
\item
  \textbf{Off-by-one in boundary} --- \texttt{lo\ \textless{}\ hi} vs \texttt{lo\ \textless{}=\ hi}; \texttt{hi\ =\ mid} vs \texttt{hi\ =\ mid\ -\ 1}
\item
  \textbf{Wrong convergence} --- \texttt{lower\_bound} needs \texttt{hi\ =\ len(arr)}, not \texttt{len(arr)\ -\ 1}
\end{itemize}

\subsection{General Interview Mistakes}\label{general-interview-mistakes}

\begin{itemize}
\tightlist
\item
  \textbf{Not clarifying constraints} --- solution for n=10 differs from n=10\^{}9
\item
  \textbf{Not discussing complexity} --- implementing a solution without analyzing it
\item
  \textbf{Jumping to code} --- not spending time designing the algorithm first
\item
  \textbf{Not testing with edge cases} --- empty input, single element, all duplicates
\item
  \textbf{Over-engineering} --- writing complex code when a simple O(n\(^2\)) is acceptable for given n
\end{itemize}

\section{Revision Cheat Sheet}\label{revision-cheat-sheet}

\subsection{Must-Know Complexities}\label{must-know-complexities}

\begin{longtable}[]{@{}
  >{\raggedright\arraybackslash}p{(\columnwidth - 4\tabcolsep) * \real{0.2973}}
  >{\raggedright\arraybackslash}p{(\columnwidth - 4\tabcolsep) * \real{0.3514}}
  >{\raggedright\arraybackslash}p{(\columnwidth - 4\tabcolsep) * \real{0.3514}}@{}}
\toprule\noalign{}
\rowcolor{acc!16}
\begin{minipage}[b]{\linewidth}\raggedright
Operation
\end{minipage} & \begin{minipage}[b]{\linewidth}\raggedright
Data Structure
\end{minipage} & \begin{minipage}[b]{\linewidth}\raggedright
Time
\end{minipage} \\
\midrule\noalign{}
\endhead
\bottomrule\noalign{}
\endlastfoot
Access by index & Array & O(1) \\
\rowcolor{zebra}
Search & Array (unsorted) & O(n) \\
Search & Array (sorted) & O(log n) \\
\rowcolor{zebra}
Insert/Delete & Array (arbitrary position) & O(n) \\
Insert/Delete & Linked List & O(1) at head; O(n) to find \\
\rowcolor{zebra}
Push/Pop & Stack/Queue & O(1) \\
Insert/Delete/Search & Hash Table & O(1) avg, O(n) worst \\
\rowcolor{zebra}
Insert/Delete/Search & BST (balanced) & O(log n) \\
Insert/Extract-Min & Min-Heap & O(log n) \\
\rowcolor{zebra}
Build heap from array & Heap & O(n) \\
DFS/BFS & Graph & O(V+E) \\
\rowcolor{zebra}
Dijkstra & Graph (non-neg) & O((V+E) log V) \\
Bellman-Ford & Graph & O(VE) \\
\rowcolor{zebra}
Floyd-Warshall & Graph & O(V\(^3\)) \\
Kruskal/Prim & Graph & O(E log E) \\
\rowcolor{zebra}
Merge Sort / Heap Sort & Array & O(n log n) \\
Quick Sort & Array & O(n log n) avg \\
\rowcolor{zebra}
Counting Sort & Array & O(n+k) \\
Binary Search & Sorted Array & O(log n) \\
\rowcolor{zebra}
Union-Find & DSU & O(\(\alpha\)(n)) amortized \\
Sieve of Eratosthenes & - & O(n log log n) \\
\rowcolor{zebra}
GCD & - & O(log(min(a,b))) \\
Fast Exponentiation & - & O(log n) \\
\end{longtable}

\subsection{Algorithm Triggers (What to Think When You See...)}\label{algorithm-triggers-what-to-think-when-you-see}

\begin{longtable}[]{@{}
  >{\raggedright\arraybackslash}p{(\columnwidth - 2\tabcolsep) * \real{0.4595}}
  >{\raggedright\arraybackslash}p{(\columnwidth - 2\tabcolsep) * \real{0.5405}}@{}}
\toprule\noalign{}
\rowcolor{acc!16}
\begin{minipage}[b]{\linewidth}\raggedright
You see...
\end{minipage} & \begin{minipage}[b]{\linewidth}\raggedright
Think...
\end{minipage} \\
\midrule\noalign{}
\endhead
\bottomrule\noalign{}
\endlastfoot
Shortest path, no weights & BFS \\
\rowcolor{zebra}
Shortest path, positive weights & Dijkstra \\
Shortest path, negative weights & Bellman-Ford \\
\rowcolor{zebra}
All-pairs shortest path & Floyd-Warshall \\
Minimum spanning tree & Kruskal (sparse) / Prim (dense) \\
\rowcolor{zebra}
"Connected components" or "groups" & Union-Find or DFS \\
"Cycle in directed graph" & Topological sort (Kahn\textquotesingle s) or DFS 3-color \\
\rowcolor{zebra}
"Dependency ordering" & Topological sort \\
"Count ways", "min/max cost" & Dynamic Programming \\
\rowcolor{zebra}
"Optimal choice at each step" & Greedy (verify with exchange argument) \\
"All possibilities / enumerate" & Backtracking \\
\rowcolor{zebra}
"Sorted array, find X" & Binary search \\
"Find minimum X such that..." & Binary search on answer \\
\rowcolor{zebra}
"Sliding window, subarray" & Two pointers or sliding window \\
"Top K elements" & Heap (size K) \\
\rowcolor{zebra}
"Anagram / frequency" & Hash map or counting array \\
"Intervals, overlapping" & Sort by start/end + greedy \\
\rowcolor{zebra}
"Parentheses, matching" & Stack \\
"Next greater element" & Monotonic stack \\
\rowcolor{zebra}
"Integer keys, small range" & Counting sort / bucket \\
"Large integers mod prime" & Modular arithmetic + Fermat \\
\rowcolor{zebra}
"Find single number in pairs" & XOR trick \\
"Power of 2 / bit patterns" & Bit manipulation \\
\rowcolor{zebra}
"String matching" & KMP / Rabin-Karp \\
"2D grid, path finding" & BFS (shortest) / DFS (any path) \\
\rowcolor{zebra}
"Subsets, combinations" & Backtracking or bitmask \\
"Repeated subproblems visible" & Memoize it \(\rightarrow\) DP \\
\end{longtable}

\subsection{Golden Rules}\label{golden-rules}

\begin{enumerate}
\def\labelenumi{\arabic{enumi}.}
\tightlist
\item
  \textbf{Always clarify constraints before coding} --- n determines your algorithm
\item
  \textbf{State complexity before implementation} --- shows algorithmic thinking
\item
  \textbf{Draw examples} --- small examples reveal patterns
\item
  \textbf{Code the brute force first if stuck} --- then optimize
\item
  \textbf{Test edge cases explicitly} --- empty, single element, all same, already sorted
\item
  \textbf{For graphs:} always check --- directed or undirected? weighted? negative edges? disconnected?
\item
  \textbf{For DP:} define state precisely, write recurrence, then code
\item
  \textbf{For greedy:} justify why local optimal = global optimal
\item
  \textbf{Binary search:} maintain invariant, never drop the answer
\item
  \textbf{Recursion:} trust the recursion --- define what it returns, use it correctly
\end{enumerate}

\clearpage
\section{Visual Reference}
\noindent A set of figures distilling the key quantitative intuitions of this topic.\par\vspace{4pt}
\visualfigure{../figures/dsa-03-algos/01-big-o-growth.pdf}{Asymptotic class, not constant factors, decides scalability. Beyond modest n, an O(n log n) algorithm crushes O(n$^2$); exponential is infeasible past ~n=30.}
\visualfigure{../figures/dsa-03-algos/02-sorting-complexity.pdf}{No single sort wins everywhere: merge sort is stable and worst-case O(n log n) but needs O(n) space; quicksort is faster in practice but O(n$^2$) worst case.}
\end{document}
